% generated by GAPDoc2LaTeX from XML source (Frank Luebeck)
\documentclass[a4paper,11pt]{report}

\usepackage[top=37mm,bottom=37mm,left=27mm,right=27mm]{geometry}
\sloppy
\pagestyle{myheadings}
\usepackage{amssymb}
\usepackage[utf8]{inputenc}
\usepackage{makeidx}
\makeindex
\usepackage{color}
\definecolor{FireBrick}{rgb}{0.5812,0.0074,0.0083}
\definecolor{RoyalBlue}{rgb}{0.0236,0.0894,0.6179}
\definecolor{RoyalGreen}{rgb}{0.0236,0.6179,0.0894}
\definecolor{RoyalRed}{rgb}{0.6179,0.0236,0.0894}
\definecolor{LightBlue}{rgb}{0.8544,0.9511,1.0000}
\definecolor{Black}{rgb}{0.0,0.0,0.0}

\definecolor{linkColor}{rgb}{0.0,0.0,0.554}
\definecolor{citeColor}{rgb}{0.0,0.0,0.554}
\definecolor{fileColor}{rgb}{0.0,0.0,0.554}
\definecolor{urlColor}{rgb}{0.0,0.0,0.554}
\definecolor{promptColor}{rgb}{0.0,0.0,0.589}
\definecolor{brkpromptColor}{rgb}{0.589,0.0,0.0}
\definecolor{gapinputColor}{rgb}{0.589,0.0,0.0}
\definecolor{gapoutputColor}{rgb}{0.0,0.0,0.0}

%%  for a long time these were red and blue by default,
%%  now black, but keep variables to overwrite
\definecolor{FuncColor}{rgb}{0.0,0.0,0.0}
%% strange name because of pdflatex bug:
\definecolor{Chapter }{rgb}{0.0,0.0,0.0}
\definecolor{DarkOlive}{rgb}{0.1047,0.2412,0.0064}


\usepackage{fancyvrb}

\usepackage{mathptmx,helvet}
\usepackage[T1]{fontenc}
\usepackage{textcomp}


\usepackage[
            pdftex=true,
            bookmarks=true,        
            a4paper=true,
            pdftitle={Written with GAPDoc},
            pdfcreator={LaTeX with hyperref package / GAPDoc},
            colorlinks=true,
            backref=page,
            breaklinks=true,
            linkcolor=linkColor,
            citecolor=citeColor,
            filecolor=fileColor,
            urlcolor=urlColor,
            pdfpagemode={UseNone}, 
           ]{hyperref}

\newcommand{\maintitlesize}{\fontsize{50}{55}\selectfont}

% write page numbers to a .pnr log file for online help
\newwrite\pagenrlog
\immediate\openout\pagenrlog =\jobname.pnr
\immediate\write\pagenrlog{PAGENRS := [}
\newcommand{\logpage}[1]{\protect\write\pagenrlog{#1, \thepage,}}
%% were never documented, give conflicts with some additional packages

\newcommand{\GAP}{\textsf{GAP}}

%% nicer description environments, allows long labels
\usepackage{enumitem}
\setdescription{style=nextline}

%% depth of toc
\setcounter{tocdepth}{1}





%% command for ColorPrompt style examples
\newcommand{\gapprompt}[1]{\color{promptColor}{\bfseries #1}}
\newcommand{\gapbrkprompt}[1]{\color{brkpromptColor}{\bfseries #1}}
\newcommand{\gapinput}[1]{\color{gapinputColor}{#1}}


\begin{document}

\logpage{[ 0, 0, 0 ]}
\begin{titlepage}
\mbox{}\vfill

\begin{center}{\maintitlesize \textbf{The \textsf{ModularGroup} Package\mbox{}}}\\
\vfill

\hypersetup{pdftitle=The \textsf{ModularGroup} Package}
\markright{\scriptsize \mbox{}\hfill The \textsf{ModularGroup} Package \hfill\mbox{}}
{\Huge \textbf{Finite\texttt{\symbol{45}}index subgroups of ${\rm (P)SL}_2({\ensuremath{\mathbb Z}})$\mbox{}}}\\
\vfill

{\Huge Version 2.0.2\mbox{}}\\[1cm]
{4 May 2026\mbox{}}\\[1cm]
\mbox{}\\[2cm]
{\Large \textbf{Sebastian Engelhardt   \mbox{}}}\\
{\Large \textbf{Luca L. Junk   \mbox{}}}\\
{\Large \textbf{Hannah Wagmann   \mbox{}}}\\
{\Large \textbf{Gabriela Weitze\texttt{\symbol{45}}Schmith{\"u}sen   \mbox{}}}\\
\hypersetup{pdfauthor=Sebastian Engelhardt   ; Luca L. Junk   ; Hannah Wagmann   ; Gabriela Weitze\texttt{\symbol{45}}Schmith{\"u}sen   }
\end{center}\vfill

\mbox{}\\
{\mbox{}\\
\small \noindent \textbf{Sebastian Engelhardt   }  Email: \href{mailto://seen00001@stud.uni-saarland.de} {\texttt{seen00001@stud.uni\texttt{\symbol{45}}saarland.de}}\\
  Homepage: \href{https://www.uni-saarland.de/lehrstuhl/weitze-schmithuesen.html} {\texttt{https://www.uni\texttt{\symbol{45}}saarland.de/lehrstuhl/weitze\texttt{\symbol{45}}schmithuesen.html}}}\\
{\mbox{}\\
\small \noindent \textbf{Luca L. Junk   }  Email: \href{mailto://junk@math.uni-sb.de} {\texttt{junk@math.uni\texttt{\symbol{45}}sb.de}}\\
  Homepage: \href{https://www.uni-saarland.de/lehrstuhl/weber-moritz/team/luca-junk.html} {\texttt{https://www.uni\texttt{\symbol{45}}saarland.de/lehrstuhl/weber\texttt{\symbol{45}}moritz/team/luca\texttt{\symbol{45}}junk.html}}}\\
{\mbox{}\\
\small \noindent \textbf{Hannah Wagmann   }  Email: \href{mailto://wagmann@math.uni-sb.de} {\texttt{wagmann@math.uni\texttt{\symbol{45}}sb.de}}\\
  Homepage: \href{https://www.uni-saarland.de/lehrstuhl/weitze-schmithuesen/team/hannah-wagmann.html} {  
            \texttt{https://www.uni-saarland.de/lehrstuhl/weitze-schmithuesen/}\\
$\indent\hspace{11.5mm}$\texttt{team/hannah-wagmann.html}
             }}\\
{\mbox{}\\
\small \noindent \textbf{Gabriela Weitze\texttt{\symbol{45}}Schmith{\"u}sen   }  Email: \href{mailto://weitze@math.uni-sb.de} {\texttt{weitze@math.uni\texttt{\symbol{45}}sb.de}}\\
  Homepage: \href{https://www.uni-saarland.de/lehrstuhl/weitze-schmithuesen/team/gabriela-weitze-schmithuesen.html} {  
            \texttt{https://www.uni-saarland.de/lehrstuhl/weitze-schmithuesen/}\\
$\indent\hspace{11.5mm}$\texttt{team/gabriela-weitze-schmithuesen.html}
             }}\\
\end{titlepage}

\newpage\setcounter{page}{2}
{\small 
\section*{Copyright}
\logpage{[ 0, 0, 1 ]}
{\copyright} 2026 by Sebastian Engelhardt, Luca L. Junk, Hannah Wagmann and
Gabriela Weitze\texttt{\symbol{45}}Schmith{\"u}sen \mbox{}}\\[1cm]
{\small 
\section*{Acknowledgements}
\logpage{[ 0, 0, 2 ]}
 Supported by Project I.8 of SFB\texttt{\symbol{45}}TRR 195 'Symbolic Tools in
Mathematics and their Application' of the German Research Foundation (DFG). \mbox{}}\\[1cm]
\newpage

\def\contentsname{Contents\logpage{[ 0, 0, 3 ]}}

\tableofcontents
\newpage

 
\chapter{\textcolor{Chapter }{Introduction}}\label{Intro}
\logpage{[ 1, 0, 0 ]}
\hyperdef{L}{X7DFB63A97E67C0A1}{}
{
  
\section{\textcolor{Chapter }{General aims of the \textsf{ModularGroup} package}}\label{IntroAims}
\logpage{[ 1, 1, 0 ]}
\hyperdef{L}{X8776EC4F822F431D}{}
{
  This \textsf{GAP} package provides methods for computing with finite\texttt{\symbol{45}}index
subgroups of the modular groups ${\rm SL}_2({\ensuremath{\mathbb Z}})$ and ${\rm PSL}_2({\ensuremath{\mathbb Z}})$. This includes, but is not limited to, computation of the generalized level,
index or cusp widths. It also implements algorithms described in \cite{hsu_1996} and \cite{hamilton_loeffler_2014} for testing if a given group is a congruence subgroup. Hence it differs from
the \href{https://gap-packages.github.io/congruence/} {Congruence} package \cite{congruence_package}, which can be used \texttt{\symbol{45}} among other things
\texttt{\symbol{45}} to construct canonical congruence subgroups of ${\rm SL}_2({\ensuremath{\mathbb Z}})$. }

 
\section{\textcolor{Chapter }{Technicalities}}\label{IntroTech}
\logpage{[ 1, 2, 0 ]}
\hyperdef{L}{X811AAE8182C9C012}{}
{
  A convenient way to represent finite\texttt{\symbol{45}}index subgroups of ${\rm SL}_2({\ensuremath{\mathbb Z}})$ is by specifying the action of generator matrices of ${\rm SL}_2({\ensuremath{\mathbb Z}})$ on the right cosets by right multiplication. For example, one could choose the
generators   
\[ S = \left( \begin{array}{rr} 0 & -1 \\ 1 & 0 \end{array} \right) \quad T =
\left( \begin{array}{rr} 1 & 1 \\ 0 & 1 \end{array} \right) \]
  and represent a subgroup as a tuple of transitive permutations $(\sigma_S,\sigma_T)$ describing the action of $S$ and $T$. This is exactly the way this package internally treats such subgroups. We
use the convention that $1$ corresponds to the coset of the identity matrix. Note that such a
representation as a tuple of permutations is only unique up to relabelling of
the cosets, i.e. up to simultaneous conjugation (fixing the $1$ coset by our convention). To not restrict \textsf{ModularGroup} to right multiplication on right cosets, getters and setters are defined in
two versions: one by right multiplication on right cosets ("\texttt{ViaRightAction}") and one by left multiplication on left cosets ("\texttt{ViaLeftAction}"). }

 }

 
\chapter{\textcolor{Chapter }{Subgroups of ${\rm SL}_2({\ensuremath{\mathbb Z}})$}}\label{SL2Z}
\logpage{[ 2, 0, 0 ]}
\hyperdef{L}{X8708256C7CDEDC88}{}
{
  For representing finite\texttt{\symbol{45}}index subgroups of ${\rm SL}_2({\ensuremath{\mathbb Z}})$, this package introduces the new object \texttt{ModularSubgroup}. As stated in the introduction, a \texttt{ModularSubgroup} essentially consists of the two permutations $\sigma_S$ and $\sigma_T$ describing the coset graph with respect to the generators $S$ and $T$ (with the convention that $1$ corresponds to the identity coset). So explicitly specifying these
permutations is the canonical way to construct a \texttt{ModularSubgroup}. \\
 Though you might not always have a coset graph of your subgroup at hand, but
rather a list of generating matrices. Therefore we implement multiple
constructors for \texttt{ModularSubgroup}: three that take as input two permutations describing the coset graph with
respect to different pairs of generators of ${\rm SL}_2({\ensuremath{\mathbb Z}})$, and one that takes a list of ${\rm SL}_2({\ensuremath{\mathbb Z}})$ matrices as generators. 
\section{\textcolor{Chapter }{Construction of modular subgroups}}\label{ModularSubgroupConstr}
\logpage{[ 2, 1, 0 ]}
\hyperdef{L}{X82B1267A8369ABC3}{}
{
  
\subsection{\textcolor{Chapter }{ModularSubgroup}}\logpage{[ 2, 1, 1 ]}
\hyperdef{L}{X7EFCADC57AA4FEC2}{}
{
\noindent\textcolor{FuncColor}{$\triangleright$\enspace\texttt{ModularSubgroupViaRightAction({\mdseries\slshape s, t})\index{ModularSubgroupViaRightAction@\texttt{ModularSubgroupViaRightAction}}
\label{ModularSubgroupViaRightAction}
}\hfill{\scriptsize (operation)}}\\
\textbf{\indent Returns:\ }
A modular subgroup.



 Constructs a \texttt{ModularSubgroup} object corresponding to the finite\texttt{\symbol{45}}index subgroup of ${\rm SL}_2({\ensuremath{\mathbb Z}})$ described by the permutations \mbox{\texttt{\mdseries\slshape s}} and \mbox{\texttt{\mdseries\slshape t}}. \\
 This constructor tests if the given permutations actually describe the coset
action of the matrices   
\[ S = \left( \begin{array}{rr} 0 & -1 \\ 1 & 0 \end{array} \right) \quad T =
\left( \begin{array}{rr} 1 & 1 \\ 0 & 1 \end{array} \right) \]
  by checking that they act transitively and satisfy the relations 
\[s^4 = (s^3 t)^3 = s^2 t s^{-2} t^{-1} = 1\]
 Upon creation, the cosets are renamed in a \href{https://www.gap-system.org/Manuals/doc/ref/chap47.html#X85B882F782D7AFD0} {standardized way} to make the internal interaction with existing \textsf{GAP} methods easier. (The fact that $1$ corresponds to the identity coset is not changed by this) 
\begin{Verbatim}[commandchars=!@|,fontsize=\small,frame=single,label=Example]
  !gapprompt@gap>| !gapinput@G := ModularSubgroupViaRightAction(|
  !gapprompt@>| !gapinput@(1,2)(3,4)(5,6)(7,8)(9,10),|
  !gapprompt@>| !gapinput@(1,4)(2,5,9,10,8)(3,7,6));|
  <modular subgroup of index 10>
\end{Verbatim}
 \noindent\textcolor{FuncColor}{$\triangleright$\enspace\texttt{ModularSubgroupViaLeftAction({\mdseries\slshape s, t})\index{ModularSubgroupViaLeftAction@\texttt{ModularSubgroupViaLeftAction}}
\label{ModularSubgroupViaLeftAction}
}\hfill{\scriptsize (operation)}}\\
\textbf{\indent Returns:\ }
A modular subgroup.



 Analogous to \texttt{ModularSubgroupViaRightAction}, but with left coset actions. }

 
\subsection{\textcolor{Chapter }{ModularSubgroupST}}\logpage{[ 2, 1, 2 ]}
\hyperdef{L}{X85492D1F7A2C97CB}{}
{
\noindent\textcolor{FuncColor}{$\triangleright$\enspace\texttt{ModularSubgroupSTViaRightAction({\mdseries\slshape s, t})\index{ModularSubgroupSTViaRightAction@\texttt{ModularSubgroupSTViaRightAction}}
\label{ModularSubgroupSTViaRightAction}
}\hfill{\scriptsize (operation)}}\\
\textbf{\indent Returns:\ }
A modular subgroup.



 Synonymous for \texttt{ModularSubgroupViaRightAction} (\ref{ModularSubgroupViaRightAction}) (see above). \noindent\textcolor{FuncColor}{$\triangleright$\enspace\texttt{ModularSubgroupSTViaLeftAction({\mdseries\slshape s, t})\index{ModularSubgroupSTViaLeftAction@\texttt{ModularSubgroupSTViaLeftAction}}
\label{ModularSubgroupSTViaLeftAction}
}\hfill{\scriptsize (operation)}}\\
\textbf{\indent Returns:\ }
A modular subgroup.



 Synonymous for \texttt{ModularSubgroupViaLeftAction} (\ref{ModularSubgroupViaLeftAction}) (see above). }

 
\subsection{\textcolor{Chapter }{ModularSubgroupRT}}\logpage{[ 2, 1, 3 ]}
\hyperdef{L}{X84D89E0B80CA1ED2}{}
{
\noindent\textcolor{FuncColor}{$\triangleright$\enspace\texttt{ModularSubgroupRTViaRightAction({\mdseries\slshape r, t})\index{ModularSubgroupRTViaRightAction@\texttt{ModularSubgroupRTViaRightAction}}
\label{ModularSubgroupRTViaRightAction}
}\hfill{\scriptsize (operation)}}\\
\textbf{\indent Returns:\ }
A modular subgroup.



 Constructs a \texttt{ModularSubgroup} object corresponding to the finite\texttt{\symbol{45}}index subgroup of ${\rm SL}_2({\ensuremath{\mathbb Z}})$ determined by the permutations \mbox{\texttt{\mdseries\slshape r}} and \mbox{\texttt{\mdseries\slshape t}} which describe the action of the matrices   
\[ R = \left( \begin{array}{rr} 1 & 0 \\ 1 & 1 \end{array} \right) \quad T =
\left( \begin{array}{rr} 1 & 1 \\ 0 & 1 \end{array} \right) \]
  on the right cosets. \\
 A check is performed if the permutations actually describe such an action on
the cosets of some subgroup. \\
 Upon creation, the cosets are renamed in a \href{https://www.gap-system.org/Manuals/doc/ref/chap47.html#X85B882F782D7AFD0} {standardized way} to make the internal interaction with existing \textsf{GAP} methods easier. (The fact that $1$ corresponds to the identity coset is not changed by this) 
\begin{Verbatim}[commandchars=!@|,fontsize=\small,frame=single,label=Example]
  !gapprompt@gap>| !gapinput@G := ModularSubgroupRTViaRightAction(|
  !gapprompt@>| !gapinput@(1,9,8,10,7)(2,6)(3,4,5),|
  !gapprompt@>| !gapinput@(1,3)(2,4,8,10,5)(6,9,7));|
  <modular subgroup of index 10>
\end{Verbatim}
 \noindent\textcolor{FuncColor}{$\triangleright$\enspace\texttt{ModularSubgroupRTViaLeftAction({\mdseries\slshape r, t})\index{ModularSubgroupRTViaLeftAction@\texttt{ModularSubgroupRTViaLeftAction}}
\label{ModularSubgroupRTViaLeftAction}
}\hfill{\scriptsize (operation)}}\\
\textbf{\indent Returns:\ }
A modular subgroup.



 Analogous to \texttt{ModularSubgroupRTViaRightAction}, but with left coset actions. }

 
\subsection{\textcolor{Chapter }{ModularSubgroupSJ}}\logpage{[ 2, 1, 4 ]}
\hyperdef{L}{X7AE9DEC97F4BC1D5}{}
{
\noindent\textcolor{FuncColor}{$\triangleright$\enspace\texttt{ModularSubgroupSJViaRightAction({\mdseries\slshape s, j})\index{ModularSubgroupSJViaRightAction@\texttt{ModularSubgroupSJViaRightAction}}
\label{ModularSubgroupSJViaRightAction}
}\hfill{\scriptsize (operation)}}\\
\textbf{\indent Returns:\ }
A modular subgroup.



 Constructs a \texttt{ModularSubgroup} object corresponding to the finite\texttt{\symbol{45}}index subgroup of ${\rm SL}_2({\ensuremath{\mathbb Z}})$ determined by the permutations \mbox{\texttt{\mdseries\slshape s}} and \mbox{\texttt{\mdseries\slshape j}} which describe the action of the matrices  
\[ S = \left( \begin{array}{rr} 0 & -1 \\ 1 & 0 \end{array} \right) \quad J =
\left( \begin{array}{rr} 0 & 1 \\ -1 & 1 \end{array} \right) \]
   on the right cosets. \\
 A check is performed if the permutations actually describe such an action on
the cosets of some subgroup. \\
 Upon creation, the cosets are renamed in a \href{https://www.gap-system.org/Manuals/doc/ref/chap47.html#X85B882F782D7AFD0} {standardized way} to make the internal interaction with existing \textsf{GAP} methods easier. (The fact that $1$ corresponds to the identity coset is not changed by this) 
\begin{Verbatim}[commandchars=!@|,fontsize=\small,frame=single,label=Example]
  !gapprompt@gap>| !gapinput@G := ModularSubgroupSJViaRightAction(|
  !gapprompt@>| !gapinput@(1,2)(3,6)(4,7)(5,9)(8,10),|
  !gapprompt@>| !gapinput@(1,5,6)(2,3,7)(4,9,10));|
  <modular subgroup of index 10>
\end{Verbatim}
 \noindent\textcolor{FuncColor}{$\triangleright$\enspace\texttt{ModularSubgrouSJViaLeftAction({\mdseries\slshape s, j})\index{ModularSubgrouSJViaLeftAction@\texttt{ModularSubgrouSJViaLeftAction}}
\label{ModularSubgrouSJViaLeftAction}
}\hfill{\scriptsize (operation)}}\\
\textbf{\indent Returns:\ }
A modular subgroup.



 Analogous to \texttt{ModularSubgroupSJViaRightAction}, but with left coset actions. }

 

\subsection{\textcolor{Chapter }{ModularSubgroup (ModularSubgroupGens)}}
\logpage{[ 2, 1, 5 ]}\nobreak
\hyperdef{L}{X7D6FE1E8831BEB59}{}
{\noindent\textcolor{FuncColor}{$\triangleright$\enspace\texttt{ModularSubgroup({\mdseries\slshape gens})\index{ModularSubgroup@\texttt{ModularSubgroup}!ModularSubgroupGens}
\label{ModularSubgroup:ModularSubgroupGens}
}\hfill{\scriptsize (operation)}}\\
\textbf{\indent Returns:\ }
A modular subgroup.



 Constructs a \texttt{ModularSubgroup} object corresponding to the finite\texttt{\symbol{45}}index subgroup of ${\rm SL}_2({\ensuremath{\mathbb Z}})$ generated by the matrices in \mbox{\texttt{\mdseries\slshape gens}}. \\
 No test is performed to check if the generated subgroup actually has finite
index! \\
 This constructor implicitly computes a coset table of the subgroup. Hence it
might be slow for very large index subgroups. 
\begin{Verbatim}[commandchars=!@|,fontsize=\small,frame=single,label=Example]
  !gapprompt@gap>| !gapinput@G := ModularSubgroup([|
  !gapprompt@>| !gapinput@[[1,2], [0,1]],|
  !gapprompt@>| !gapinput@[[1,0], [2,1]],|
  !gapprompt@>| !gapinput@[[-1,0], [0,-1]]|
  !gapprompt@>| !gapinput@]);|
  <modular subgroup of index 6>
\end{Verbatim}
 }

 }

 
\section{\textcolor{Chapter }{Getters for the coset action}}\logpage{[ 2, 2, 0 ]}
\hyperdef{L}{X7A38ED567DB2C5EC}{}
{
  
\subsection{\textcolor{Chapter }{SAction}}\logpage{[ 2, 2, 1 ]}
\hyperdef{L}{X7806F137806CE532}{}
{
\noindent\textcolor{FuncColor}{$\triangleright$\enspace\texttt{SRightAction({\mdseries\slshape G})\index{SRightAction@\texttt{SRightAction}}
\label{SRightAction}
}\hfill{\scriptsize (operation)}}\\
\textbf{\indent Returns:\ }
A permutation.



 Returns the permutation $\sigma_S$ describing the action of the matrix $S$ on the right cosets of \mbox{\texttt{\mdseries\slshape G}}. \noindent\textcolor{FuncColor}{$\triangleright$\enspace\texttt{SLeftAction({\mdseries\slshape G})\index{SLeftAction@\texttt{SLeftAction}}
\label{SLeftAction}
}\hfill{\scriptsize (operation)}}\\
\textbf{\indent Returns:\ }
A permutation.



 Analogous to \texttt{SRightAction}, but with left coset actions. }

 
\subsection{\textcolor{Chapter }{TAction}}\logpage{[ 2, 2, 2 ]}
\hyperdef{L}{X845C82DF798AAC68}{}
{
\noindent\textcolor{FuncColor}{$\triangleright$\enspace\texttt{TRightAction({\mdseries\slshape G})\index{TRightAction@\texttt{TRightAction}}
\label{TRightAction}
}\hfill{\scriptsize (operation)}}\\
\textbf{\indent Returns:\ }
A permutation.



 Returns the permutation $\sigma_T$ describing the action of the matrix $T$ on the right cosets of \mbox{\texttt{\mdseries\slshape G}}. \noindent\textcolor{FuncColor}{$\triangleright$\enspace\texttt{TLeftAction({\mdseries\slshape G})\index{TLeftAction@\texttt{TLeftAction}}
\label{TLeftAction}
}\hfill{\scriptsize (operation)}}\\
\textbf{\indent Returns:\ }
A permutation.



 Analogous to \texttt{TRightAction}, but with left coset actions. }

 
\subsection{\textcolor{Chapter }{RAction}}\logpage{[ 2, 2, 3 ]}
\hyperdef{L}{X8261818C871C963B}{}
{
\noindent\textcolor{FuncColor}{$\triangleright$\enspace\texttt{RRightAction({\mdseries\slshape G})\index{RRightAction@\texttt{RRightAction}}
\label{RRightAction}
}\hfill{\scriptsize (operation)}}\\
\textbf{\indent Returns:\ }
A permutation.



 Returns the permutation $\sigma_R$ describing the action of the matrix $R$ on the right cosets of \mbox{\texttt{\mdseries\slshape G}}. \noindent\textcolor{FuncColor}{$\triangleright$\enspace\texttt{RLeftAction({\mdseries\slshape G})\index{RLeftAction@\texttt{RLeftAction}}
\label{RLeftAction}
}\hfill{\scriptsize (operation)}}\\
\textbf{\indent Returns:\ }
A permutation.



 Analogous to \texttt{RRightAction}, but with left coset actions. }

 
\subsection{\textcolor{Chapter }{JAction}}\logpage{[ 2, 2, 4 ]}
\hyperdef{L}{X87229CA6862A5FBE}{}
{
\noindent\textcolor{FuncColor}{$\triangleright$\enspace\texttt{JRightAction({\mdseries\slshape G})\index{JRightAction@\texttt{JRightAction}}
\label{JRightAction}
}\hfill{\scriptsize (operation)}}\\
\textbf{\indent Returns:\ }
A permutation.



 Returns the permutation $\sigma_J$ describing the action of the matrix $J$ on the right cosets of \mbox{\texttt{\mdseries\slshape G}}. \noindent\textcolor{FuncColor}{$\triangleright$\enspace\texttt{JLeftAction({\mdseries\slshape G})\index{JLeftAction@\texttt{JLeftAction}}
\label{JLeftAction}
}\hfill{\scriptsize (operation)}}\\
\textbf{\indent Returns:\ }
A permutation.



 Analogous to \texttt{JRightAction}, but with left coset actions. }

 
\subsection{\textcolor{Chapter }{CosetActionOf}}\logpage{[ 2, 2, 5 ]}
\hyperdef{L}{X859D10B779F28DFD}{}
{
\noindent\textcolor{FuncColor}{$\triangleright$\enspace\texttt{CosetRightActionOf({\mdseries\slshape A, G})\index{CosetRightActionOf@\texttt{CosetRightActionOf}}
\label{CosetRightActionOf}
}\hfill{\scriptsize (operation)}}\\
\textbf{\indent Returns:\ }
A permutation.



 Returns the permutation $\sigma_A$ describing the action of the matrix $A \in {\rm SL}_2({\ensuremath{\mathbb Z}})$ on the right cosets of \mbox{\texttt{\mdseries\slshape G}}. \noindent\textcolor{FuncColor}{$\triangleright$\enspace\texttt{CosetLeftActionOf({\mdseries\slshape A, G})\index{CosetLeftActionOf@\texttt{CosetLeftActionOf}}
\label{CosetLeftActionOf}
}\hfill{\scriptsize (operation)}}\\
\textbf{\indent Returns:\ }
A permutation.



 Analogous to \texttt{CosetRightActionOf}, but with left coset actions. }

 }

 
\section{\textcolor{Chapter }{Computing with modular subgroups}}\label{ModularSubgroupAlgs}
\logpage{[ 2, 3, 0 ]}
\hyperdef{L}{X7B2F1C71867A167F}{}
{
  

\subsection{\textcolor{Chapter }{Index (IndexSL2Z)}}
\logpage{[ 2, 3, 1 ]}\nobreak
\hyperdef{L}{X789D6D848365E376}{}
{\noindent\textcolor{FuncColor}{$\triangleright$\enspace\texttt{Index({\mdseries\slshape G})\index{Index@\texttt{Index}!IndexSL2Z}
\label{Index:IndexSL2Z}
}\hfill{\scriptsize (attribute)}}\\
\textbf{\indent Returns:\ }
A natural number.



 For a given modular subgroup \mbox{\texttt{\mdseries\slshape G}} this method returns its index in ${\rm SL}_2({\ensuremath{\mathbb Z}})$. As \mbox{\texttt{\mdseries\slshape G}} is internally stored as permutations $(s,t)$ this is just \texttt{LargestMovedPoint([s,t])} (or $1$ if the permutations are trivial). }

 
\begin{Verbatim}[commandchars=!@|,fontsize=\small,frame=single,label=Example]
  !gapprompt@gap>| !gapinput@G := ModularSubgroup((1,2)(3,5)(4,6), (1,3)(2,4)(5,6));|
  <modular subgroup of index 6>
  !gapprompt@gap>| !gapinput@Index(G);|
  6
\end{Verbatim}
 

\subsection{\textcolor{Chapter }{GeneralizedLevel (GeneralizedLevelSL2Z)}}
\logpage{[ 2, 3, 2 ]}\nobreak
\hyperdef{L}{X81C43D8E7FB0EB1C}{}
{\noindent\textcolor{FuncColor}{$\triangleright$\enspace\texttt{GeneralizedLevel({\mdseries\slshape G})\index{GeneralizedLevel@\texttt{GeneralizedLevel}!GeneralizedLevelSL2Z}
\label{GeneralizedLevel:GeneralizedLevelSL2Z}
}\hfill{\scriptsize (attribute)}}\\
\textbf{\indent Returns:\ }
A natural number.



 This method calculates the general Wohlfahrt level (i.e. the lowest common
multiple of all cusp widths) of \mbox{\texttt{\mdseries\slshape G}} as defined in \cite{wohlfahrt1964}. }

 
\begin{Verbatim}[commandchars=!@|,fontsize=\small,frame=single,label=Example]
  !gapprompt@gap>| !gapinput@G := ModularSubgroup((1,2)(3,5)(4,6), (1,3)(2,4)(5,6));|
  <modular subgroup of index 6>
  !gapprompt@gap>| !gapinput@GeneralizedLevel(G);|
  2
\end{Verbatim}
 
\subsection{\textcolor{Chapter }{CosetRepresentatives}}\logpage{[ 2, 3, 3 ]}
\hyperdef{L}{X806A331E8509C9F4}{}
{
\noindent\textcolor{FuncColor}{$\triangleright$\enspace\texttt{RightCosetRepresentatives({\mdseries\slshape G})\index{RightCosetRepresentatives@\texttt{RightCosetRepresentatives}!RightCosetRepresentativesSL2Z}
\label{RightCosetRepresentatives:RightCosetRepresentativesSL2Z}
}\hfill{\scriptsize (attribute)}}\\
\textbf{\indent Returns:\ }
A list of words.



 This function returns a list of representatives of the (right) cosets of \mbox{\texttt{\mdseries\slshape G}} as words in $S$ and $T$. 
\begin{Verbatim}[commandchars=!@|,fontsize=\small,frame=single,label=Example]
  !gapprompt@gap>| !gapinput@G := ModularSubgroup((1,2),(2,3));|
  <modular subgroup of index 3>
  !gapprompt@gap>| !gapinput@RightCosetRepresentatives(G);|
  [ <identity ...>, S, S*T ]
\end{Verbatim}
 \noindent\textcolor{FuncColor}{$\triangleright$\enspace\texttt{LeftCosetRepresentatives({\mdseries\slshape G})\index{LeftCosetRepresentatives@\texttt{LeftCosetRepresentatives}!LeftCosetRepresentativesSL2Z}
\label{LeftCosetRepresentatives:LeftCosetRepresentativesSL2Z}
}\hfill{\scriptsize (attribute)}}\\
\textbf{\indent Returns:\ }
A list of words.



 This function returns a list of representatives of the (left) cosets of \mbox{\texttt{\mdseries\slshape G}} as words in $S$ and $T$. }

 

\subsection{\textcolor{Chapter }{WordGeneratorsOfGroup (GeneratorsOfGroupSL2Z)}}
\logpage{[ 2, 3, 4 ]}\nobreak
\hyperdef{L}{X878C488F84FD9039}{}
{\noindent\textcolor{FuncColor}{$\triangleright$\enspace\texttt{WordGeneratorsOfGroup({\mdseries\slshape G})\index{WordGeneratorsOfGroup@\texttt{WordGeneratorsOfGroup}!GeneratorsOfGroupSL2Z}
\label{WordGeneratorsOfGroup:GeneratorsOfGroupSL2Z}
}\hfill{\scriptsize (attribute)}}\\
\textbf{\indent Returns:\ }
A list of words.



 Calculates a list of generators (as words in $S$ and $T$) of \mbox{\texttt{\mdseries\slshape G}}. This list might include redundant generators (or even duplicates). }

 
\begin{Verbatim}[commandchars=!@|,fontsize=\small,frame=single,label=Example]
  !gapprompt@gap>| !gapinput@G := ModularSubgroup((1,2)(3,5)(4,6), (1,3)(2,4)(5,6));|
  <modular subgroup of index 6>
  !gapprompt@gap>| !gapinput@WordGeneratorsOfGroup(G);|
  [ S^-2, T^-2, S*T^-2*S^-1 ]
\end{Verbatim}
 

\subsection{\textcolor{Chapter }{GeneratorsOfGroup}}
\logpage{[ 2, 3, 5 ]}\nobreak
\hyperdef{L}{X79C44528864044C5}{}
{\noindent\textcolor{FuncColor}{$\triangleright$\enspace\texttt{GeneratorsOfGroup({\mdseries\slshape G})\index{GeneratorsOfGroup@\texttt{GeneratorsOfGroup}}
\label{GeneratorsOfGroup}
}\hfill{\scriptsize (attribute)}}\\
\textbf{\indent Returns:\ }
A list of matrices.



 Calculates a list of generator matrices of \mbox{\texttt{\mdseries\slshape G}}. This list might include redundant generators (or even duplicates). }

 
\begin{Verbatim}[commandchars=!@|,fontsize=\small,frame=single,label=Example]
  !gapprompt@gap>| !gapinput@G := ModularSubgroup((1,2)(3,5)(4,6), (1,3)(2,4)(5,6));|
  <modular subgroup of index 6>
  !gapprompt@gap>| !gapinput@GeneratorsOfGroup(G);|
  [ [ [ -1, 0 ], [ 0, -1 ] ], [ [ 1, -2 ], [ 0, 1 ] ], [ [ 1, 0 ], [ 2, 1 ] ] ]
\end{Verbatim}
 

\subsection{\textcolor{Chapter }{IsCongruence (IsCongruenceSL2Z)}}
\logpage{[ 2, 3, 6 ]}\nobreak
\hyperdef{L}{X873013BD876F7DAD}{}
{\noindent\textcolor{FuncColor}{$\triangleright$\enspace\texttt{IsCongruence({\mdseries\slshape G})\index{IsCongruence@\texttt{IsCongruence}!IsCongruenceSL2Z}
\label{IsCongruence:IsCongruenceSL2Z}
}\hfill{\scriptsize (attribute)}}\\
\textbf{\indent Returns:\ }
True or false.



 This method test whether a given modular subgroup \mbox{\texttt{\mdseries\slshape G}} is a congruence subgroup. It is essentially an implementation of an algorithm
described in \cite{hamilton_loeffler_2014}. }

 
\begin{Verbatim}[commandchars=!@|,fontsize=\small,frame=single,label=Example]
  !gapprompt@gap>| !gapinput@G := ModularSubgroup([|
  !gapprompt@>| !gapinput@[[1,2],[0,1]],|
  !gapprompt@>| !gapinput@[[1,0],[2,1]]|
  !gapprompt@>| !gapinput@]);|
  <modular subgroup of index 12>
  !gapprompt@gap>| !gapinput@IsCongruence(G);|
  true
\end{Verbatim}
 

\subsection{\textcolor{Chapter }{Cusps (CuspsSL2Z)}}
\logpage{[ 2, 3, 7 ]}\nobreak
\hyperdef{L}{X78BF0B867D965A6C}{}
{\noindent\textcolor{FuncColor}{$\triangleright$\enspace\texttt{Cusps({\mdseries\slshape G})\index{Cusps@\texttt{Cusps}!CuspsSL2Z}
\label{Cusps:CuspsSL2Z}
}\hfill{\scriptsize (attribute)}}\\
\textbf{\indent Returns:\ }
A list of rational numbers and infinity.



 This method computes a list of inequivalent cusp representatives with respect
to \mbox{\texttt{\mdseries\slshape G}}. }

 
\begin{Verbatim}[commandchars=!@|,fontsize=\small,frame=single,label=Example]
  !gapprompt@gap>| !gapinput@G := ModularSubgroup(|
  !gapprompt@>| !gapinput@(1,2)(3,6)(4,8)(5,9)(7,11)(10,13)(12,15)(14,17)(16,19)(18,21)(20,23)(22,24),|
  !gapprompt@>| !gapinput@(1,3,7,4)(2,5)(6,9,8,12,14,10)(11,13,16,20,18,15)(17,21,22,19)(23,24)|
  !gapprompt@>| !gapinput@);|
  <modular subgroup of index 24>
  !gapprompt@gap>| !gapinput@Cusps(G);|
  [ infinity, 0, 1, 2, 3/2, 5/3 ]
\end{Verbatim}
 

\subsection{\textcolor{Chapter }{CuspWidth (CuspWidthSL2Z)}}
\logpage{[ 2, 3, 8 ]}\nobreak
\hyperdef{L}{X83E84049842EC55A}{}
{\noindent\textcolor{FuncColor}{$\triangleright$\enspace\texttt{CuspWidth({\mdseries\slshape c, G})\index{CuspWidth@\texttt{CuspWidth}!CuspWidthSL2Z}
\label{CuspWidth:CuspWidthSL2Z}
}\hfill{\scriptsize (operation)}}\\
\textbf{\indent Returns:\ }
A natural number.



 This method takes as input a cusp \mbox{\texttt{\mdseries\slshape c}} (a rational number or infinity) and a modular group \mbox{\texttt{\mdseries\slshape G}} and calculates the width of this cusp with respect to \mbox{\texttt{\mdseries\slshape G}}. }

 
\begin{Verbatim}[commandchars=!@|,fontsize=\small,frame=single,label=Example]
  !gapprompt@gap>| !gapinput@G := ModularSubgroup(|
  !gapprompt@>| !gapinput@(1,2,6,3)(4,11,15,12)(5,13,16,14)(7,17,9,18)(8,19,10,20)(21,24,22,23),|
  !gapprompt@>| !gapinput@(1,4,5)(2,7,8)(3,9,10)(6,15,16)(11,20,21)(12,19,22)(13,23,17)(14,24,18)|
  !gapprompt@>| !gapinput@);|
  <modular subgroup of index 24>
  !gapprompt@gap>| !gapinput@CuspWidth(-1, G);|
  3
  !gapprompt@gap>| !gapinput@CuspWidth(infinity, G);|
  3
\end{Verbatim}
 

\subsection{\textcolor{Chapter }{CuspsEquivalent (CuspsEquivalentSL2Z)}}
\logpage{[ 2, 3, 9 ]}\nobreak
\hyperdef{L}{X80EBA8E17B6DC9F1}{}
{\noindent\textcolor{FuncColor}{$\triangleright$\enspace\texttt{CuspsEquivalent({\mdseries\slshape p, q, G})\index{CuspsEquivalent@\texttt{CuspsEquivalent}!CuspsEquivalentSL2Z}
\label{CuspsEquivalent:CuspsEquivalentSL2Z}
}\hfill{\scriptsize (operation)}}\\
\textbf{\indent Returns:\ }
True or false.



 Takes two cusps \mbox{\texttt{\mdseries\slshape p}} and \mbox{\texttt{\mdseries\slshape q}} and a modular subgroup \mbox{\texttt{\mdseries\slshape G}} and checks if they are equivalent modulo \mbox{\texttt{\mdseries\slshape G}}, i.e. if there exists a matrix $A \in G$ with $Ap = q$. }

 
\begin{Verbatim}[commandchars=!@|,fontsize=\small,frame=single,label=Example]
  !gapprompt@gap>| !gapinput@G := ModularSubgroup(|
  !gapprompt@>| !gapinput@(1,2,6,3)(4,11,15,12)(5,13,16,14)(7,17,9,18)(8,19,10,20)(21,24,22,23),|
  !gapprompt@>| !gapinput@(1,4,5)(2,7,8)(3,9,10)(6,15,16)(11,20,21)(12,19,22)(13,23,17)(14,24,18)|
  !gapprompt@>| !gapinput@);|
  <modular subgroup of index 24>
  !gapprompt@gap>| !gapinput@CuspsEquivalent(infinity, 1, G);|
  false
  !gapprompt@gap>| !gapinput@CuspsEquivalent(-1, 1/2, G);|
  true
\end{Verbatim}
 

\subsection{\textcolor{Chapter }{CosetRepresentativeOfCusp (CosetRepresentativeOfCuspSL2Z)}}
\logpage{[ 2, 3, 10 ]}\nobreak
\hyperdef{L}{X7A73C80F7AA9FCF1}{}
{\noindent\textcolor{FuncColor}{$\triangleright$\enspace\texttt{CosetRepresentativeOfCusp({\mdseries\slshape c, G})\index{CosetRepresentativeOfCusp@\texttt{CosetRepresentativeOfCusp}!CosetRepresentativeOfCuspSL2Z}
\label{CosetRepresentativeOfCusp:CosetRepresentativeOfCuspSL2Z}
}\hfill{\scriptsize (operation)}}\\
\textbf{\indent Returns:\ }
A word in S and T.



 For a cusp \mbox{\texttt{\mdseries\slshape c}} this function returns a right coset representative $A$ of \mbox{\texttt{\mdseries\slshape G}} such that $A \infty $ and \mbox{\texttt{\mdseries\slshape c}} are equivalent with respect to \mbox{\texttt{\mdseries\slshape G}}. }

 
\begin{Verbatim}[commandchars=!@|,fontsize=\small,frame=single,label=Example]
  !gapprompt@gap>| !gapinput@G := ModularSubgroup(|
  !gapprompt@>| !gapinput@(1,2,6,3)(4,11,15,12)(5,13,16,14)(7,17,9,18)(8,19,10,20)(21,24,22,23),|
  !gapprompt@>| !gapinput@(1,4,5)(2,7,8)(3,9,10)(6,15,16)(11,20,21)(12,19,22)(13,23,17)(14,24,18)|
  !gapprompt@>| !gapinput@);|
  <modular subgroup of index 24>
  !gapprompt@gap>| !gapinput@CosetRepresentativeOfCusp(4, G);|
  T*S
\end{Verbatim}
 

\subsection{\textcolor{Chapter }{IndexModN (IndexModNSL2Z)}}
\logpage{[ 2, 3, 11 ]}\nobreak
\hyperdef{L}{X7FEB5D3B853CA831}{}
{\noindent\textcolor{FuncColor}{$\triangleright$\enspace\texttt{IndexModN({\mdseries\slshape G, N})\index{IndexModN@\texttt{IndexModN}!IndexModNSL2Z}
\label{IndexModN:IndexModNSL2Z}
}\hfill{\scriptsize (operation)}}\\
\textbf{\indent Returns:\ }
A natural number.



 For a modular subgroup \mbox{\texttt{\mdseries\slshape G}} and a natural number \mbox{\texttt{\mdseries\slshape N}} this method calculates the index of the projection $\bar{G}$ of $G$ in ${\rm SL}_2({\ensuremath{\mathbb Z}}/N{\ensuremath{\mathbb Z}})$. }

 
\begin{Verbatim}[commandchars=!@|,fontsize=\small,frame=single,label=Example]
  !gapprompt@gap>| !gapinput@G := ModularSubgroup((1,2)(3,5)(4,6), (1,3)(2,4)(5,6));|
  <modular subgroup of index 6>
  !gapprompt@gap>| !gapinput@IndexModN(G, 2);|
  6
\end{Verbatim}
 

\subsection{\textcolor{Chapter }{Deficiency (DeficiencySL2Z)}}
\logpage{[ 2, 3, 12 ]}\nobreak
\hyperdef{L}{X7E6754897B892B4A}{}
{\noindent\textcolor{FuncColor}{$\triangleright$\enspace\texttt{Deficiency({\mdseries\slshape G, N})\index{Deficiency@\texttt{Deficiency}!DeficiencySL2Z}
\label{Deficiency:DeficiencySL2Z}
}\hfill{\scriptsize (operation)}}\\
\textbf{\indent Returns:\ }
A natural number.



 For a modular subgroup \mbox{\texttt{\mdseries\slshape G}} and a natural number \mbox{\texttt{\mdseries\slshape N}} this method calculates the so\texttt{\symbol{45}}called \emph{deficiency} of \mbox{\texttt{\mdseries\slshape G}} from being a congruence subgroup of level \mbox{\texttt{\mdseries\slshape N}}. \\
 The deficiency of a finite\texttt{\symbol{45}}index subgroup $\Gamma$ of ${\rm SL}_2({\ensuremath{\mathbb Z}})$ was introduced in \cite{weitze_schmithuesen_deficiency}. It is defined as the index $[\Gamma(N) \colon \Gamma(N) \cap \Gamma]$ where $\Gamma(N)$ is the principal congruence subgroup of level $N$. }

 
\begin{Verbatim}[commandchars=!@|,fontsize=\small,frame=single,label=Example]
  !gapprompt@gap>| !gapinput@G := ModularSubgroup([|
  !gapprompt@>| !gapinput@[[1,2],[0,1]],|
  !gapprompt@>| !gapinput@[[1,0],[2,1]]|
  !gapprompt@>| !gapinput@]);|
  <modular subgroup of index 12>
  !gapprompt@gap>| !gapinput@Deficiency(G, 2);|
  2
  !gapprompt@gap>| !gapinput@Deficiency(G, 4);|
  1
\end{Verbatim}
 

\subsection{\textcolor{Chapter }{Deficiency (DeficiencySL2ZAttr)}}
\logpage{[ 2, 3, 13 ]}\nobreak
\hyperdef{L}{X7B704B49846C54E7}{}
{\noindent\textcolor{FuncColor}{$\triangleright$\enspace\texttt{Deficiency({\mdseries\slshape G})\index{Deficiency@\texttt{Deficiency}!DeficiencySL2ZAttr}
\label{Deficiency:DeficiencySL2ZAttr}
}\hfill{\scriptsize (attribute)}}\\
\textbf{\indent Returns:\ }
A natural number.



 Shorthand for \texttt{Deficiency(G, GeneralizedLevel(G))}. }

 
\begin{Verbatim}[commandchars=!@|,fontsize=\small,frame=single,label=Example]
  !gapprompt@gap>| !gapinput@G := ModularSubgroup([|
  !gapprompt@>| !gapinput@[[1,2],[0,1]],|
  !gapprompt@>| !gapinput@[[1,0],[2,1]]|
  !gapprompt@>| !gapinput@]);|
  <modular subgroup of index 12>
  !gapprompt@gap>| !gapinput@Deficiency(G);|
  1
  !gapprompt@gap>| !gapinput@Deficiency(G, GeneralizedLevel(G));|
  1
\end{Verbatim}
 

\subsection{\textcolor{Chapter }{Projectivization}}
\logpage{[ 2, 3, 14 ]}\nobreak
\hyperdef{L}{X81E057F386298DB2}{}
{\noindent\textcolor{FuncColor}{$\triangleright$\enspace\texttt{Projectivization({\mdseries\slshape G})\index{Projectivization@\texttt{Projectivization}}
\label{Projectivization}
}\hfill{\scriptsize (operation)}}\\
\textbf{\indent Returns:\ }
A projective modular subgroup.



 For a given modular subgroup \mbox{\texttt{\mdseries\slshape G}} this function calculates its image $\bar{\mbox{\texttt{\mdseries\slshape G}}}$ under the projection $\pi \colon {\rm SL}_2({\ensuremath{\mathbb Z}}) \rightarrow {\rm
PSL}_2({\ensuremath{\mathbb Z}})$. }

 
\begin{Verbatim}[commandchars=!@|,fontsize=\small,frame=single,label=Example]
  !gapprompt@gap>| !gapinput@G := ModularSubgroup([|
  !gapprompt@>| !gapinput@[[1,2],[0,1]],|
  !gapprompt@>| !gapinput@[[1,0],[2,1]]|
  !gapprompt@>| !gapinput@]);|
  <modular subgroup of index 12>
  !gapprompt@gap>| !gapinput@Projectivization(G);|
  <projective modular subgroup of index 6>
\end{Verbatim}
 

\subsection{\textcolor{Chapter }{ConjugateGroup (ConjugateGroupSL2Z)}}
\logpage{[ 2, 3, 15 ]}\nobreak
\hyperdef{L}{X7AD35F397834117F}{}
{\noindent\textcolor{FuncColor}{$\triangleright$\enspace\texttt{ConjugateGroup({\mdseries\slshape G, A})\index{ConjugateGroup@\texttt{ConjugateGroup}!ConjugateGroupSL2Z}
\label{ConjugateGroup:ConjugateGroupSL2Z}
}\hfill{\scriptsize (operation)}}\\
\textbf{\indent Returns:\ }
A ModularSubgroup.



 Conjugates the group \mbox{\texttt{\mdseries\slshape G}} by the ${\rm SL}_2({\ensuremath{\mathbb Z}})$ matrix \mbox{\texttt{\mdseries\slshape A}} and returns the group $A^{-1}*G*A$. }

 

\subsection{\textcolor{Chapter }{NormalCore (NormalCoreSL2Z)}}
\logpage{[ 2, 3, 16 ]}\nobreak
\hyperdef{L}{X874DC4FA8436FCBC}{}
{\noindent\textcolor{FuncColor}{$\triangleright$\enspace\texttt{NormalCore({\mdseries\slshape G})\index{NormalCore@\texttt{NormalCore}!NormalCoreSL2Z}
\label{NormalCore:NormalCoreSL2Z}
}\hfill{\scriptsize (attribute)}}\\
\textbf{\indent Returns:\ }
A modular subgroup.



 Calculates the normal core of \mbox{\texttt{\mdseries\slshape G}} in ${\rm SL}_2({\ensuremath{\mathbb Z}})$, i.e. the maximal subgroup of \mbox{\texttt{\mdseries\slshape G}} that is normal in ${\rm SL}_2({\ensuremath{\mathbb Z}})$. }

 
\begin{Verbatim}[commandchars=!@|,fontsize=\small,frame=single,label=Example]
  !gapprompt@gap>| !gapinput@G := ModularSubgroup([|
  !gapprompt@>| !gapinput@[[1,2],[0,1]],|
  !gapprompt@>| !gapinput@[[1,0],[2,1]]|
  !gapprompt@>| !gapinput@]);|
  <modular subgroup of index 12>
  !gapprompt@gap>| !gapinput@NormalCore(G);|
  <modular subgroup of index 48>
\end{Verbatim}
 

\subsection{\textcolor{Chapter }{QuotientByNormalCore (QuotientByNormalCoreSL2Z)}}
\logpage{[ 2, 3, 17 ]}\nobreak
\hyperdef{L}{X7D645ABF7BF4600B}{}
{\noindent\textcolor{FuncColor}{$\triangleright$\enspace\texttt{QuotientByNormalCore({\mdseries\slshape G})\index{QuotientByNormalCore@\texttt{QuotientByNormalCore}!QuotientByNormalCoreSL2Z}
\label{QuotientByNormalCore:QuotientByNormalCoreSL2Z}
}\hfill{\scriptsize (attribute)}}\\
\textbf{\indent Returns:\ }
A finite group.



 Calculates the quotient of ${\rm SL}_2({\ensuremath{\mathbb Z}})$ by the normal core of \mbox{\texttt{\mdseries\slshape G}}. }

 
\begin{Verbatim}[commandchars=!@|,fontsize=\small,frame=single,label=Example]
  !gapprompt@gap>| !gapinput@G := ModularSubgroup([|
  !gapprompt@>| !gapinput@[[1,2],[0,1]],|
  !gapprompt@>| !gapinput@[[1,0],[2,1]]|
  !gapprompt@>| !gapinput@]);|
  <modular subgroup of index 12>
  !gapprompt@gap>| !gapinput@QuotientByNormalCore(G);|
  <permutation group with 2 generators>
\end{Verbatim}
 

\subsection{\textcolor{Chapter }{AssociatedCharacterTable (AssociatedCharacterTableSL2Z)}}
\logpage{[ 2, 3, 18 ]}\nobreak
\hyperdef{L}{X874C4C4584FB5D03}{}
{\noindent\textcolor{FuncColor}{$\triangleright$\enspace\texttt{AssociatedCharacterTable({\mdseries\slshape G})\index{AssociatedCharacterTable@\texttt{AssociatedCharacterTable}!AssociatedCharacterTableSL2Z}
\label{AssociatedCharacterTable:AssociatedCharacterTableSL2Z}
}\hfill{\scriptsize (attribute)}}\\
\textbf{\indent Returns:\ }
A character table.



 Returns the character table of ${\rm SL}_2({\ensuremath{\mathbb Z}})/N$ where $N$ is the normal core of \mbox{\texttt{\mdseries\slshape G}}. }

 
\begin{Verbatim}[commandchars=!@|,fontsize=\small,frame=single,label=Example]
  !gapprompt@gap>| !gapinput@G := ModularSubgroup([|
  !gapprompt@>| !gapinput@[[1,2],[0,1]],|
  !gapprompt@>| !gapinput@[[1,0],[2,1]]|
  !gapprompt@>| !gapinput@]);|
  <modular subgroup of index 12>
  !gapprompt@gap>| !gapinput@AssociatedCharacterTable(G);|
  CharacterTable( <permutation group of size 48 with 2 generators> )
\end{Verbatim}
 

\subsection{\textcolor{Chapter }{IsElementOf (IsElementOfSL2Z)}}
\logpage{[ 2, 3, 19 ]}\nobreak
\hyperdef{L}{X79E8C4D97BF50B74}{}
{\noindent\textcolor{FuncColor}{$\triangleright$\enspace\texttt{IsElementOf({\mdseries\slshape A, G})\index{IsElementOf@\texttt{IsElementOf}!IsElementOfSL2Z}
\label{IsElementOf:IsElementOfSL2Z}
}\hfill{\scriptsize (operation)}}\\
\textbf{\indent Returns:\ }
True or false.



 This function checks if a given matrix \mbox{\texttt{\mdseries\slshape A}} is an element of the modular subgroup \mbox{\texttt{\mdseries\slshape G}}. }

 
\begin{Verbatim}[commandchars=!@|,fontsize=\small,frame=single,label=Example]
  !gapprompt@gap>| !gapinput@G := ModularSubgroup([|
  !gapprompt@>| !gapinput@[[1,2],[0,1]],|
  !gapprompt@>| !gapinput@[[1,0],[2,1]]|
  !gapprompt@>| !gapinput@]);|
  <modular subgroup of index 12>
  !gapprompt@gap>| !gapinput@IsElementOf([[-1,0],[0,-1]], G);|
  false
  !gapprompt@gap>| !gapinput@IsElementOf([[1,4],[0,1]], G);|
  true
\end{Verbatim}
 

\subsection{\textcolor{Chapter }{IsWordElementOf (IsWordElementOfSL2ZString)}}
\logpage{[ 2, 3, 20 ]}\nobreak
\hyperdef{L}{X801BE17A799532FB}{}
{\noindent\textcolor{FuncColor}{$\triangleright$\enspace\texttt{IsWordElementOf({\mdseries\slshape s, G})\index{IsWordElementOf@\texttt{IsWordElementOf}!IsWordElementOfSL2ZString}
\label{IsWordElementOf:IsWordElementOfSL2ZString}
}\hfill{\scriptsize (operation)}}\\
\textbf{\indent Returns:\ }
True or false.



 This function checks if a given string \mbox{\texttt{\mdseries\slshape s}} represents an element of the modular subgroup \mbox{\texttt{\mdseries\slshape G}}, written in the generators \mbox{\texttt{\mdseries\slshape S}} and \mbox{\texttt{\mdseries\slshape T}}. }

 
\begin{Verbatim}[commandchars=!@|,fontsize=\small,frame=single,label=Example]
  !gapprompt@gap>| !gapinput@G := ModularSubgroup([|
  !gapprompt@>| !gapinput@[[1,2],[0,1]],|
  !gapprompt@>| !gapinput@[[1,0],[2,1]]|
  !gapprompt@>| !gapinput@]);|
  <modular subgroup of index 12>
  !gapprompt@gap>| !gapinput@IsWordElementOf("S^4 * T^2", G);|
  true
\end{Verbatim}
 

\subsection{\textcolor{Chapter }{IsWordElementOf (IsWordElementOfSL2Z)}}
\logpage{[ 2, 3, 21 ]}\nobreak
\hyperdef{L}{X7EFB6B6C80ACBE93}{}
{\noindent\textcolor{FuncColor}{$\triangleright$\enspace\texttt{IsWordElementOf({\mdseries\slshape w, G})\index{IsWordElementOf@\texttt{IsWordElementOf}!IsWordElementOfSL2Z}
\label{IsWordElementOf:IsWordElementOfSL2Z}
}\hfill{\scriptsize (operation)}}\\
\textbf{\indent Returns:\ }
True or false.



 This function checks if a given word \mbox{\texttt{\mdseries\slshape w}} in the generators \mbox{\texttt{\mdseries\slshape S}} and \mbox{\texttt{\mdseries\slshape T}} represents an element of the modular subgroup \mbox{\texttt{\mdseries\slshape G}}. }

 
\begin{Verbatim}[commandchars=!@|,fontsize=\small,frame=single,label=Example]
  !gapprompt@gap>| !gapinput@G := ModularSubgroup([|
  !gapprompt@>| !gapinput@[[1,2],[0,1]],|
  !gapprompt@>| !gapinput@[[1,0],[2,1]]|
  !gapprompt@>| !gapinput@]);|
  <modular subgroup of index 12>
  !gapprompt@gap>| !gapinput@F := FreeGroup("S", "T");|
  <free group on the generators [ S, T ]>
  !gapprompt@gap>| !gapinput@S := F.1;|
  S
  !gapprompt@gap>| !gapinput@T := F.2;|
  T
  !gapprompt@gap>| !gapinput@SL2Z := F / ParseRelators([S, T], "S^4, (S^3*T)^3, S^2*T*S^-2*T^-1");|
  <fp group on the generators [ S, T ]>
  !gapprompt@gap>| !gapinput@S := GeneratorsOfGroup(SL2Z)[1];|
  S
  !gapprompt@gap>| !gapinput@T := GeneratorsOfGroup(SL2Z)[2];|
  T
  !gapprompt@gap>| !gapinput@IsWordElementOf(S^4 * T^2, G);|
  true
\end{Verbatim}
 

\subsection{\textcolor{Chapter }{Genus (GenusSL2Z)}}
\logpage{[ 2, 3, 22 ]}\nobreak
\hyperdef{L}{X87B255C4784B8CBC}{}
{\noindent\textcolor{FuncColor}{$\triangleright$\enspace\texttt{Genus({\mdseries\slshape G})\index{Genus@\texttt{Genus}!GenusSL2Z}
\label{Genus:GenusSL2Z}
}\hfill{\scriptsize (attribute)}}\\
\textbf{\indent Returns:\ }
A non\texttt{\symbol{45}}negative integer.



 Computes the genus of the quotient $G \setminus {\ensuremath{\mathbb H}}$ via an algorithm described in \cite{schmithuesen_alg}. }

 
\begin{Verbatim}[commandchars=!@|,fontsize=\small,frame=single,label=Example]
  !gapprompt@gap>| !gapinput@G := ModularSubgroup((1,2),(2,3));|
  <modular subgroup of index 3>
  !gapprompt@gap>| !gapinput@Genus(G);|
  0
\end{Verbatim}
 }

 
\section{\textcolor{Chapter }{Miscellaneous}}\label{ModularSubgroupMisc}
\logpage{[ 2, 4, 0 ]}
\hyperdef{L}{X7C5563A37D566DA5}{}
{
  The following functions are mostly helper functions used internally and are
only documented for sake of completeness. 

\subsection{\textcolor{Chapter }{DefinesCosetActionST}}
\logpage{[ 2, 4, 1 ]}\nobreak
\hyperdef{L}{X7DF96C8184DCF28E}{}
{\noindent\textcolor{FuncColor}{$\triangleright$\enspace\texttt{DefinesCosetActionST({\mdseries\slshape s, t})\index{DefinesCosetActionST@\texttt{DefinesCosetActionST}}
\label{DefinesCosetActionST}
}\hfill{\scriptsize (operation)}}\\
\textbf{\indent Returns:\ }
True or false.



 Checks if two given permutations \mbox{\texttt{\mdseries\slshape s}} and \mbox{\texttt{\mdseries\slshape t}} describe the action of the generator matrices $S$ and $T$ on the cosets of some subgroup. This is the case if they satisfy the relations 
\[s^4 = (s^3 t)^3 = s^2 t s^{-2} t^{-1} = 1\]
 and act transitively. }

 
\begin{Verbatim}[commandchars=!@|,fontsize=\small,frame=single,label=Example]
  !gapprompt@gap>| !gapinput@s := (1,2)(3,4)(5,6)(7,8)(9,10);;|
  !gapprompt@gap>| !gapinput@t := (1,4)(2,5,9,10,8)(3,7,6);;|
  !gapprompt@gap>| !gapinput@DefinesCosetActionST(s,t);|
  true
\end{Verbatim}
 

\subsection{\textcolor{Chapter }{DefinesCosetActionRT}}
\logpage{[ 2, 4, 2 ]}\nobreak
\hyperdef{L}{X7C68DF9579A2DA8B}{}
{\noindent\textcolor{FuncColor}{$\triangleright$\enspace\texttt{DefinesCosetActionRT({\mdseries\slshape r, t})\index{DefinesCosetActionRT@\texttt{DefinesCosetActionRT}}
\label{DefinesCosetActionRT}
}\hfill{\scriptsize (operation)}}\\
\textbf{\indent Returns:\ }
True or false.



 Checks if two given permutations \mbox{\texttt{\mdseries\slshape r}} and \mbox{\texttt{\mdseries\slshape t}} describe the action of the generator matrices $R$ and $T$ on the cosets of some subgroup. This is the case if they satisfy the relations 
\[(r t^{-1} r)^4 = ((r t^{-1} r)^3 t)^3 = (r t^{-1} r)^2 t (r t^{-1} r)^{-2}
t^{-1} = 1\]
 and act transitively. }

 
\begin{Verbatim}[commandchars=!@|,fontsize=\small,frame=single,label=Example]
  !gapprompt@gap>| !gapinput@r := (1,9,8,10,7)(2,6)(3,4,5);;|
  !gapprompt@gap>| !gapinput@t := (1,3)(2,4,8,10,5)(6,9,7);;|
  !gapprompt@gap>| !gapinput@DefinesCosetActionRT(r,t);|
  true
\end{Verbatim}
 

\subsection{\textcolor{Chapter }{DefinesCosetActionSJ}}
\logpage{[ 2, 4, 3 ]}\nobreak
\hyperdef{L}{X82599F577B4F1C8A}{}
{\noindent\textcolor{FuncColor}{$\triangleright$\enspace\texttt{DefinesCosetActionSJ({\mdseries\slshape s, j})\index{DefinesCosetActionSJ@\texttt{DefinesCosetActionSJ}}
\label{DefinesCosetActionSJ}
}\hfill{\scriptsize (operation)}}\\
\textbf{\indent Returns:\ }
True or false.



 Checks if two given permutations \mbox{\texttt{\mdseries\slshape s}} and \mbox{\texttt{\mdseries\slshape j}} describe the action of the generator matrices $S$ and $J$ on the cosets of some subgroup. This is the case if they satisfy the relations 
\[s^4 = (s^3 j^{-1} s^{-1})^3 = s^2 j^{-1} s^{-2} j = 1\]
 and act transitively. }

 
\begin{Verbatim}[commandchars=!@|,fontsize=\small,frame=single,label=Example]
  !gapprompt@gap>| !gapinput@s := (1,2)(3,4)(5,6)(7,8)(9,10);;|
  !gapprompt@gap>| !gapinput@j := (1,5,6)(2,3,7)(4,9,10);;|
  !gapprompt@gap>| !gapinput@DefinesCosetActionSJ(s,j);|
  true
\end{Verbatim}
 
\subsection{\textcolor{Chapter }{CosetActionFromGenerators}}\logpage{[ 2, 4, 4 ]}
\hyperdef{L}{X7D96A4DE84701927}{}
{
\noindent\textcolor{FuncColor}{$\triangleright$\enspace\texttt{RightCosetActionFromGenerators({\mdseries\slshape gens})\index{RightCosetActionFromGenerators@\texttt{RightCosetActionFromGenerators}}
\label{RightCosetActionFromGenerators}
}\hfill{\scriptsize (operation)}}\\
\textbf{\indent Returns:\ }
A tuple of permutations.



 Takes a list of generator matrices and calculates the right coset graph (as
two permutations $\sigma_S$ and $\sigma_T$) of the generated subgroup of ${\rm SL}_2({\ensuremath{\mathbb Z}})$. 
\begin{Verbatim}[commandchars=!@|,fontsize=\small,frame=single,label=Example]
  !gapprompt@gap>| !gapinput@RightCosetActionFromGenerators([|
  !gapprompt@>| !gapinput@[[1,2],[0,1]],|
  !gapprompt@>| !gapinput@[[1,0],[2,1]]|
  !gapprompt@>| !gapinput@]);|
  [ (1,2,5,3)(4,8,10,9)(6,11,7,12), (1,4)(2,6)(3,7)(5,10)(8,12,9,11) ]
\end{Verbatim}
 \noindent\textcolor{FuncColor}{$\triangleright$\enspace\texttt{LeftCosetActionFromGenerators({\mdseries\slshape gens})\index{LeftCosetActionFromGenerators@\texttt{LeftCosetActionFromGenerators}}
\label{LeftCosetActionFromGenerators}
}\hfill{\scriptsize (operation)}}\\
\textbf{\indent Returns:\ }
A tuple of permutations.



 Analogous to \texttt{RightCosetActionFromGenerators}, but with left coset actions. }

 

\subsection{\textcolor{Chapter }{STDecomposition}}
\logpage{[ 2, 4, 5 ]}\nobreak
\hyperdef{L}{X7CC5EB727FF9442E}{}
{\noindent\textcolor{FuncColor}{$\triangleright$\enspace\texttt{STDecomposition({\mdseries\slshape A})\index{STDecomposition@\texttt{STDecomposition}}
\label{STDecomposition}
}\hfill{\scriptsize (operation)}}\\
\textbf{\indent Returns:\ }
A word in $S$ and $T$.



 Takes a matrix $\mbox{\texttt{\mdseries\slshape A}} \in {\rm SL}_2({\ensuremath{\mathbb Z}})$ and decomposes it into a word in the generator matrices $S$ and $T$. }

 
\begin{Verbatim}[commandchars=!@|,fontsize=\small,frame=single,label=Example]
  !gapprompt@gap>| !gapinput@M := [ [ 4, 3 ], [ -3, -2 ] ];;|
  !gapprompt@gap>| !gapinput@STDecomposition(M);|
  S^2*T^-1*S^-1*T^2*S^-1*T^-1*S^-1
\end{Verbatim}
 

\subsection{\textcolor{Chapter }{RTDecomposition}}
\logpage{[ 2, 4, 6 ]}\nobreak
\hyperdef{L}{X875A56CC78893727}{}
{\noindent\textcolor{FuncColor}{$\triangleright$\enspace\texttt{RTDecomposition({\mdseries\slshape A})\index{RTDecomposition@\texttt{RTDecomposition}}
\label{RTDecomposition}
}\hfill{\scriptsize (operation)}}\\
\textbf{\indent Returns:\ }
A word in $R$ and $T$.



 Takes a matrix $\mbox{\texttt{\mdseries\slshape A}} \in {\rm SL}_2({\ensuremath{\mathbb Z}})$ and decomposes it into a word in the generator matrices $R$ and $T$. }

 
\begin{Verbatim}[commandchars=!@|,fontsize=\small,frame=single,label=Example]
  !gapprompt@gap>| !gapinput@M := [ [ 4, 3 ], [ -3, -2 ] ];;|
  !gapprompt@gap>| !gapinput@RTDecomposition(M);|
  (R*T^-1*R)^2*T^-1*R^-1*(T*R^-1*T)^2*R^-1*T^-1*R^-1*T*R^-1
\end{Verbatim}
 

\subsection{\textcolor{Chapter }{SJDecomposition}}
\logpage{[ 2, 4, 7 ]}\nobreak
\hyperdef{L}{X86177DEC82BD57D3}{}
{\noindent\textcolor{FuncColor}{$\triangleright$\enspace\texttt{SJDecomposition({\mdseries\slshape A})\index{SJDecomposition@\texttt{SJDecomposition}}
\label{SJDecomposition}
}\hfill{\scriptsize (operation)}}\\
\textbf{\indent Returns:\ }
A word in $S$ and $J$.



 Takes a matrix $\mbox{\texttt{\mdseries\slshape A}} \in {\rm SL}_2({\ensuremath{\mathbb Z}})$ and decomposes it into a word in the generator matrices $S$ and $J$. }

 
\begin{Verbatim}[commandchars=!@|,fontsize=\small,frame=single,label=Example]
  !gapprompt@gap>| !gapinput@M := [ [ 4, 3 ], [ -3, -2 ] ];;|
  !gapprompt@gap>| !gapinput@SJDecomposition(M);|
  S^3*J*(S^-1*J^-1)^2*S^-1*J*S^-1
\end{Verbatim}
 

\subsection{\textcolor{Chapter }{STDecompositionAsList}}
\logpage{[ 2, 4, 8 ]}\nobreak
\hyperdef{L}{X7AFE18407D135E5D}{}
{\noindent\textcolor{FuncColor}{$\triangleright$\enspace\texttt{STDecompositionAsList({\mdseries\slshape A})\index{STDecompositionAsList@\texttt{STDecompositionAsList}}
\label{STDecompositionAsList}
}\hfill{\scriptsize (operation)}}\\
\textbf{\indent Returns:\ }
A list representing a word in $S$ and $T$.



 Takes a matrix $\mbox{\texttt{\mdseries\slshape A}} \in {\rm SL}_2({\ensuremath{\mathbb Z}})$ and decomposes it into a word in the generator matrices $S$ and $T$. The word is represented as a list in the format [[generator, exponent], ...
] }

 
\begin{Verbatim}[commandchars=!@|,fontsize=\small,frame=single,label=Example]
  !gapprompt@gap>| !gapinput@M := [ [ 4, 3 ], [ -3, -2 ] ];;|
  !gapprompt@gap>| !gapinput@STDecompositionAsList(M);|
  [ [ "S", 2 ], [ "T", -1 ], [ "S", -1 ], [ "T", 2 ], [ "S", -1 ], [ "T", -1 ],
    [ "S", -1 ], [ "T", 0 ] ]
\end{Verbatim}
 }

 }

 
\chapter{\textcolor{Chapter }{Subgroups of ${\rm PSL}_2({\ensuremath{\mathbb Z}})$}}\label{PSL2Z}
\logpage{[ 3, 0, 0 ]}
\hyperdef{L}{X7F6A6DE77A993D2B}{}
{
  Analogous to finite\texttt{\symbol{45}}index subgroups of ${\rm SL}_2({\ensuremath{\mathbb Z}})$, we define a new type \texttt{ProjectiveModularSubgroup} for representing subgroups of ${\rm PSL}_2({\ensuremath{\mathbb Z}})$. It consists essentially of two permutations $\sigma_{\overline{S}}$ and $\sigma_{\overline{T}}$ describing the action of $\overline{S}$ and $\overline{T}$ on the cosets of the given subgroup, where $\overline{S}$ and $\overline{T}$ are the images of $S$ and $T$ in ${\rm PSL}_2({\ensuremath{\mathbb Z}})$. \\
 The methods implemented for ${\rm PSL}_2({\ensuremath{\mathbb Z}})$ subgroups are mostly the same as for ${\rm SL}_2({\ensuremath{\mathbb Z}})$ subgroups and behave more or less identically. Nevertheless we list them here. 
\section{\textcolor{Chapter }{Construction of projective modular subgroups}}\label{ProjectiveModularSubgroupConstr}
\logpage{[ 3, 1, 0 ]}
\hyperdef{L}{X809E1406842D5ADC}{}
{
  
\subsection{\textcolor{Chapter }{ProjectiveModularSubgroup}}\logpage{[ 3, 1, 1 ]}
\hyperdef{L}{X8253E7267976838F}{}
{
\noindent\textcolor{FuncColor}{$\triangleright$\enspace\texttt{ProjectiveModularSubgroupViaRightAction({\mdseries\slshape s, t})\index{ProjectiveModularSubgroupViaRightAction@\texttt{Projective}\-\texttt{Modular}\-\texttt{Subgroup}\-\texttt{Via}\-\texttt{Right}\-\texttt{Action}}
\label{ProjectiveModularSubgroupViaRightAction}
}\hfill{\scriptsize (operation)}}\\
\textbf{\indent Returns:\ }
A projective modular subgroup.



 Constructs a \texttt{ProjectiveModularSubgroup} object corresponding to the finite\texttt{\symbol{45}}index subgroup of ${\rm PSL}_2({\ensuremath{\mathbb Z}})$ described by the permutations \mbox{\texttt{\mdseries\slshape s}} and \mbox{\texttt{\mdseries\slshape t}}. \\
 This constructor tests if the given permutations actually describe the coset
action of $\overline{S}$ and $\overline{T}$ on some subgroup by checking that they act transitively and satisfy the
relations 
\[s^2 = (s t)^3 = 1\]
 Upon creation, the cosets are renamed in a \href{https://www.gap-system.org/Manuals/doc/ref/chap47.html#X85B882F782D7AFD0} {standardized way} to make the internal interaction with extisting \textsf{GAP} methods easier. 
\begin{Verbatim}[commandchars=!@|,fontsize=\small,frame=single,label=Example]
  !gapprompt@gap>| !gapinput@G := ProjectiveModularSubgroupViaRightAction(|
  !gapprompt@>| !gapinput@(1,2)(3,4)(5,6)(7,8)(9,10),|
  !gapprompt@>| !gapinput@(1,4)(2,5,9,10,8)(3,7,6));|
  <projective modular subgroup of index 10>
\end{Verbatim}
 If you want to construct a \texttt{ProjectiveModularSubgroup} from a list of generators, you can lift each generator to a matrix in ${\rm SL}_2({\ensuremath{\mathbb Z}})$, construct from these a \texttt{ModularSubgroup}, and then project it to ${\rm PSL}_2({\ensuremath{\mathbb Z}})$ via \texttt{Projectivization}. \noindent\textcolor{FuncColor}{$\triangleright$\enspace\texttt{ProjectiveModularSubgroupViaLeftAction({\mdseries\slshape s, t})\index{ProjectiveModularSubgroupViaLeftAction@\texttt{Projective}\-\texttt{Modular}\-\texttt{Subgroup}\-\texttt{Via}\-\texttt{Left}\-\texttt{Action}}
\label{ProjectiveModularSubgroupViaLeftAction}
}\hfill{\scriptsize (operation)}}\\
\textbf{\indent Returns:\ }
A projective modular subgroup.



 Analogous to \texttt{ProjectiveModularSubgroupViaRightAction}, but with left coset actions. }

 }

 
\section{\textcolor{Chapter }{Getters for the coset action}}\logpage{[ 3, 2, 0 ]}
\hyperdef{L}{X7A38ED567DB2C5EC}{}
{
  
\subsection{\textcolor{Chapter }{SAction}}\logpage{[ 3, 2, 1 ]}
\hyperdef{L}{X7806F137806CE532}{}
{
\noindent\textcolor{FuncColor}{$\triangleright$\enspace\texttt{SRightAction({\mdseries\slshape G})\index{SRightAction@\texttt{SRightAction}!SRightActionPSL2Z}
\label{SRightAction:SRightActionPSL2Z}
}\hfill{\scriptsize (operation)}}\\
\textbf{\indent Returns:\ }
A permutation.



 Returns the permutation $\sigma_{\overline{S}}$ describing the action of $\overline{S}$ on the right cosets of \mbox{\texttt{\mdseries\slshape G}}. \\
 \noindent\textcolor{FuncColor}{$\triangleright$\enspace\texttt{SLeftAction({\mdseries\slshape G})\index{SLeftAction@\texttt{SLeftAction}!SLeftActionPSL2Z}
\label{SLeftAction:SLeftActionPSL2Z}
}\hfill{\scriptsize (operation)}}\\
\textbf{\indent Returns:\ }
A permutation.



 Analogous to \texttt{SRightAction}, but with left coset actions. }

 
\subsection{\textcolor{Chapter }{TAction}}\logpage{[ 3, 2, 2 ]}
\hyperdef{L}{X845C82DF798AAC68}{}
{
\noindent\textcolor{FuncColor}{$\triangleright$\enspace\texttt{TRightAction({\mdseries\slshape G})\index{TRightAction@\texttt{TRightAction}!TRightActionPSL2Z}
\label{TRightAction:TRightActionPSL2Z}
}\hfill{\scriptsize (operation)}}\\
\textbf{\indent Returns:\ }
A permutation.



 Returns the permutation $\sigma_{\overline{T}}$ describing the action of $\overline{T}$ on the right cosets of \mbox{\texttt{\mdseries\slshape G}}. \noindent\textcolor{FuncColor}{$\triangleright$\enspace\texttt{TLeftAction({\mdseries\slshape G})\index{TLeftAction@\texttt{TLeftAction}!TLeftActionPSL2Z}
\label{TLeftAction:TLeftActionPSL2Z}
}\hfill{\scriptsize (operation)}}\\
\textbf{\indent Returns:\ }
A permutation.



 Analogous to \texttt{TRightAction}, but with left coset actions. }

 }

 
\section{\textcolor{Chapter }{Computing with projective modular subgroups}}\label{ProjectiveModularSubgroupAlgs}
\logpage{[ 3, 3, 0 ]}
\hyperdef{L}{X7942A9F0799620E7}{}
{
  

\subsection{\textcolor{Chapter }{Index (IndexPSL2Z)}}
\logpage{[ 3, 3, 1 ]}\nobreak
\hyperdef{L}{X8400F2D48474810D}{}
{\noindent\textcolor{FuncColor}{$\triangleright$\enspace\texttt{Index({\mdseries\slshape G})\index{Index@\texttt{Index}!IndexPSL2Z}
\label{Index:IndexPSL2Z}
}\hfill{\scriptsize (attribute)}}\\
\textbf{\indent Returns:\ }
A natural number.



 For a given projective modular subgroup \mbox{\texttt{\mdseries\slshape G}} this method returns its index in ${\rm PSL}_2({\ensuremath{\mathbb Z}})$. As \mbox{\texttt{\mdseries\slshape G}} is internally stored as permutations $(s,t)$ this is just 
\begin{verbatim}  
        LargestMovedPoint(s,t)
      
\end{verbatim}
 (or $1$ if the permutations are trivial). }

 
\begin{Verbatim}[commandchars=!@|,fontsize=\small,frame=single,label=Example]
  !gapprompt@gap>| !gapinput@G := ProjectiveModularSubgroup((1,2),(2,3));|
  <projective modular subgroup of index 3>
  !gapprompt@gap>| !gapinput@Index(G);|
  3
\end{Verbatim}
 

\subsection{\textcolor{Chapter }{GeneralizedLevel (GeneralizedLevelPSL2Z)}}
\logpage{[ 3, 3, 2 ]}\nobreak
\hyperdef{L}{X80F454A187CF784D}{}
{\noindent\textcolor{FuncColor}{$\triangleright$\enspace\texttt{GeneralizedLevel({\mdseries\slshape G})\index{GeneralizedLevel@\texttt{GeneralizedLevel}!GeneralizedLevelPSL2Z}
\label{GeneralizedLevel:GeneralizedLevelPSL2Z}
}\hfill{\scriptsize (attribute)}}\\
\textbf{\indent Returns:\ }
A natural number.



 This method calculates the general Wohlfahrt level (i.e. the lowest common
multiple of all cusp widths) of \mbox{\texttt{\mdseries\slshape G}} as defined in \cite{wohlfahrt1964}. }

 
\begin{Verbatim}[commandchars=!@|,fontsize=\small,frame=single,label=Example]
  !gapprompt@gap>| !gapinput@G := ProjectiveModularSubgroup(|
  !gapprompt@>| !gapinput@(1,2)(3,6)(4,8)(5,9)(7,11)(10,13)(12,15)(14,17)(16,19)(18,21)(20,23)(22,24),|
  !gapprompt@>| !gapinput@(1,3,7,4)(2,5)(6,9,8,12,14,10)(11,13,16,20,18,15)(17,21,22,19)(23,24)|
  !gapprompt@>| !gapinput@);|
  <projective modular subgroup of index 24>
  !gapprompt@gap>| !gapinput@GeneralizedLevel(G);|
  12
\end{Verbatim}
 
\subsection{\textcolor{Chapter }{CosetRepresentatives}}\logpage{[ 3, 3, 3 ]}
\hyperdef{L}{X806A331E8509C9F4}{}
{
\noindent\textcolor{FuncColor}{$\triangleright$\enspace\texttt{RightCosetRepresentatives({\mdseries\slshape G})\index{RightCosetRepresentatives@\texttt{RightCosetRepresentatives}!RightCosetRepresentativesPSL2Z}
\label{RightCosetRepresentatives:RightCosetRepresentativesPSL2Z}
}\hfill{\scriptsize (attribute)}}\\
\textbf{\indent Returns:\ }
A list of words.



 This function returns a list of representatives of the (right) cosets of \mbox{\texttt{\mdseries\slshape G}} as words in $\overline{S}$ and $\overline{T}$. 
\begin{Verbatim}[commandchars=!@|,fontsize=\small,frame=single,label=Example]
  !gapprompt@gap>| !gapinput@G := ProjectiveModularSubgroup((1,2),(2,3));|
  <projective modular subgroup of index 3>
  !gapprompt@gap>| !gapinput@RightCosetRepresentatives(G);|
  [ <identity ...>, S, S*T ]
\end{Verbatim}
 \noindent\textcolor{FuncColor}{$\triangleright$\enspace\texttt{LeftCosetRepresentatives({\mdseries\slshape G})\index{LeftCosetRepresentatives@\texttt{LeftCosetRepresentatives}!LeftCosetRepresentativesPSL2Z}
\label{LeftCosetRepresentatives:LeftCosetRepresentativesPSL2Z}
}\hfill{\scriptsize (attribute)}}\\
\textbf{\indent Returns:\ }
A list of words.



 Analogous to \texttt{LeftCosetRepresentatives}, but with left coset actions. }

 

\subsection{\textcolor{Chapter }{WordGeneratorsOfGroup (GeneratorsOfGroupPSL2Z)}}
\logpage{[ 3, 3, 4 ]}\nobreak
\hyperdef{L}{X78B9F8C97BC4E107}{}
{\noindent\textcolor{FuncColor}{$\triangleright$\enspace\texttt{WordGeneratorsOfGroup({\mdseries\slshape G})\index{WordGeneratorsOfGroup@\texttt{WordGeneratorsOfGroup}!GeneratorsOfGroupPSL2Z}
\label{WordGeneratorsOfGroup:GeneratorsOfGroupPSL2Z}
}\hfill{\scriptsize (attribute)}}\\
\textbf{\indent Returns:\ }
A list of words.



 Calculates a list of generators (as words in $\overline{S}$ and $\overline{T}$) of \mbox{\texttt{\mdseries\slshape G}}. This list might include redundant generators. }

 
\begin{Verbatim}[commandchars=!@|,fontsize=\small,frame=single,label=Example]
  !gapprompt@gap>| !gapinput@G := ProjectiveModularSubgroup((1,2)(3,5)(4,6), (1,3)(2,4)(5,6));|
  <projective modular subgroup of index 6>
  !gapprompt@gap>| !gapinput@WordGeneratorsOfGroup(G);|
  [ T^-2, S*T^-2*S^-1 ]
\end{Verbatim}
 

\subsection{\textcolor{Chapter }{IsCongruence (IsCongruencePSL2Z)}}
\logpage{[ 3, 3, 5 ]}\nobreak
\hyperdef{L}{X7B94C25C81A424DF}{}
{\noindent\textcolor{FuncColor}{$\triangleright$\enspace\texttt{IsCongruence({\mdseries\slshape G})\index{IsCongruence@\texttt{IsCongruence}!IsCongruencePSL2Z}
\label{IsCongruence:IsCongruencePSL2Z}
}\hfill{\scriptsize (attribute)}}\\
\textbf{\indent Returns:\ }
True or false.



 This method test whether a given modular subgroup \mbox{\texttt{\mdseries\slshape G}} is a congruence subgroup. It is essentially an implementation of an algorithm
described in \cite{hsu_1996}. }

 
\begin{Verbatim}[commandchars=!@|,fontsize=\small,frame=single,label=Example]
  !gapprompt@gap>| !gapinput@G := ProjectiveModularSubgroup(|
  !gapprompt@>| !gapinput@(1,2)(3,5)(4,6),|
  !gapprompt@>| !gapinput@(1,3)(2,4)(5,6)|
  !gapprompt@>| !gapinput@);|
  <projective modular subgroup of index 6>
  !gapprompt@gap>| !gapinput@IsCongruence(G);|
  true
\end{Verbatim}
 

\subsection{\textcolor{Chapter }{Cusps (CuspsPSL2Z)}}
\logpage{[ 3, 3, 6 ]}\nobreak
\hyperdef{L}{X7E507BE47A873817}{}
{\noindent\textcolor{FuncColor}{$\triangleright$\enspace\texttt{Cusps({\mdseries\slshape G})\index{Cusps@\texttt{Cusps}!CuspsPSL2Z}
\label{Cusps:CuspsPSL2Z}
}\hfill{\scriptsize (attribute)}}\\
\textbf{\indent Returns:\ }
A list of rational numbers and infinity.



 This method computes a list of inequivalent cusp representatives with respect
to \mbox{\texttt{\mdseries\slshape G}}. }

 
\begin{Verbatim}[commandchars=!@|,fontsize=\small,frame=single,label=Example]
  !gapprompt@gap>| !gapinput@G := ProjectiveModularSubgroup(|
  !gapprompt@>| !gapinput@(1,2)(3,6)(4,8)(5,9)(7,11)(10,13)(12,15)(14,17)(16,19)(18,21)(20,23)(22,24),|
  !gapprompt@>| !gapinput@(1,3,7,4)(2,5)(6,9,8,12,14,10)(11,13,16,20,18,15)(17,21,22,19)(23,24)|
  !gapprompt@>| !gapinput@);|
  <projective modular subgroup of index 24>
  !gapprompt@gap>| !gapinput@Cusps(G);|
  [ infinity, 0, 1, 2, 3/2, 5/3 ]
\end{Verbatim}
 

\subsection{\textcolor{Chapter }{CuspWidth (CuspWidthPSL2Z)}}
\logpage{[ 3, 3, 7 ]}\nobreak
\hyperdef{L}{X821D38048561DBE4}{}
{\noindent\textcolor{FuncColor}{$\triangleright$\enspace\texttt{CuspWidth({\mdseries\slshape c, G})\index{CuspWidth@\texttt{CuspWidth}!CuspWidthPSL2Z}
\label{CuspWidth:CuspWidthPSL2Z}
}\hfill{\scriptsize (operation)}}\\
\textbf{\indent Returns:\ }
A natural number.



 This method takes as input a cusp \mbox{\texttt{\mdseries\slshape c}} (a rational number or infinity) and a modular group \mbox{\texttt{\mdseries\slshape G}} and calculates the width of this cusp with respect to \mbox{\texttt{\mdseries\slshape G}}. }

 
\begin{Verbatim}[commandchars=!@|,fontsize=\small,frame=single,label=Example]
  !gapprompt@gap>| !gapinput@G := ProjectiveModularSubgroup(|
  !gapprompt@>| !gapinput@(1,2)(3,7)(4,8)(5,9)(6,10)(11,12),|
  !gapprompt@>| !gapinput@(1,3,4)(2,5,6)(7,10,11)(8,12,9)|
  !gapprompt@>| !gapinput@);|
  <projective modular subgroup of index 12>
  !gapprompt@gap>| !gapinput@CuspWidth(-1, G);|
  3
  !gapprompt@gap>| !gapinput@CuspWidth(infinity, G);|
  3
\end{Verbatim}
 

\subsection{\textcolor{Chapter }{CuspsEquivalent (CuspsEquivalentPSL2Z)}}
\logpage{[ 3, 3, 8 ]}\nobreak
\hyperdef{L}{X7C280B5D7F4FF4CC}{}
{\noindent\textcolor{FuncColor}{$\triangleright$\enspace\texttt{CuspsEquivalent({\mdseries\slshape p, q, G})\index{CuspsEquivalent@\texttt{CuspsEquivalent}!CuspsEquivalentPSL2Z}
\label{CuspsEquivalent:CuspsEquivalentPSL2Z}
}\hfill{\scriptsize (operation)}}\\
\textbf{\indent Returns:\ }
True or false.



 Takes two cusps \mbox{\texttt{\mdseries\slshape p}} and \mbox{\texttt{\mdseries\slshape q}} and a projective modular subgroup \mbox{\texttt{\mdseries\slshape G}} and checks if they are equivalent modulo \mbox{\texttt{\mdseries\slshape G}}, i.e. if there exists $A \in G$ with $Ap = q$. }

 
\begin{Verbatim}[commandchars=!@|,fontsize=\small,frame=single,label=Example]
  !gapprompt@gap>| !gapinput@G := ProjectiveModularSubgroup(|
  !gapprompt@>| !gapinput@(1,2)(3,7)(4,8)(5,9)(6,10)(11,12),|
  !gapprompt@>| !gapinput@(1,3,4)(2,5,6)(7,10,11)(8,12,9)|
  !gapprompt@>| !gapinput@);|
  <projective modular subgroup of index 12>
  !gapprompt@gap>| !gapinput@CuspsEquivalent(infinity, 1, G);|
  false
  !gapprompt@gap>| !gapinput@CuspsEquivalent(-1, 1/2, G);|
  true
\end{Verbatim}
 

\subsection{\textcolor{Chapter }{CosetRepresentativeOfCusp (CosetRepresentativeOfCuspPSL2Z)}}
\logpage{[ 3, 3, 9 ]}\nobreak
\hyperdef{L}{X7F29CCC382880A6F}{}
{\noindent\textcolor{FuncColor}{$\triangleright$\enspace\texttt{CosetRepresentativeOfCusp({\mdseries\slshape c, G})\index{CosetRepresentativeOfCusp@\texttt{CosetRepresentativeOfCusp}!CosetRepresentativeOfCuspPSL2Z}
\label{CosetRepresentativeOfCusp:CosetRepresentativeOfCuspPSL2Z}
}\hfill{\scriptsize (operation)}}\\
\textbf{\indent Returns:\ }
A word in S and T.



 For a cusp \mbox{\texttt{\mdseries\slshape c}} this function returns a right coset representative $A$ of \mbox{\texttt{\mdseries\slshape G}} such that $A \infty $ and \mbox{\texttt{\mdseries\slshape c}} are equivalent with respect to \mbox{\texttt{\mdseries\slshape G}}. }

 
\begin{Verbatim}[commandchars=!@|,fontsize=\small,frame=single,label=Example]
  !gapprompt@gap>| !gapinput@G := ProjectiveModularSubgroup(|
  !gapprompt@>| !gapinput@(1,2)(3,7)(4,8)(5,9)(6,10)(11,12),|
  !gapprompt@>| !gapinput@(1,3,4)(2,5,6)(7,10,11)(8,12,9)|
  !gapprompt@>| !gapinput@);|
  <projective modular subgroup of index 12>
  !gapprompt@gap>| !gapinput@CosetRepresentativeOfCusp(4, G);|
  T*S
\end{Verbatim}
 

\subsection{\textcolor{Chapter }{LiftToSL2ZEven}}
\logpage{[ 3, 3, 10 ]}\nobreak
\hyperdef{L}{X83153565839269C2}{}
{\noindent\textcolor{FuncColor}{$\triangleright$\enspace\texttt{LiftToSL2ZEven({\mdseries\slshape G})\index{LiftToSL2ZEven@\texttt{LiftToSL2ZEven}}
\label{LiftToSL2ZEven}
}\hfill{\scriptsize (operation)}}\\
\textbf{\indent Returns:\ }
A modular subgroup.



 Lifts a given subgroup \mbox{\texttt{\mdseries\slshape G}} of ${\rm PSL}_2({\ensuremath{\mathbb Z}})$ to an even subgroup of ${\rm SL}_2({\ensuremath{\mathbb Z}})$, i.e. a group that contains $-1$ and whose projection to ${\rm PSL}_2({\ensuremath{\mathbb Z}})$ is \mbox{\texttt{\mdseries\slshape G}}. }

 
\begin{Verbatim}[commandchars=!@|,fontsize=\small,frame=single,label=Example]
  !gapprompt@gap>| !gapinput@G := ProjectiveModularSubgroup(|
  !gapprompt@>| !gapinput@(1,2)(3,7)(4,8)(5,9)(6,10)(11,12),|
  !gapprompt@>| !gapinput@(1,3,4)(2,5,6)(7,10,11)(8,12,9)|
  !gapprompt@>| !gapinput@);|
  <projective modular subgroup of index 12>
  !gapprompt@gap>| !gapinput@LiftToSL2ZEven(G);|
  <modular subgroup of index 12>
\end{Verbatim}
 

\subsection{\textcolor{Chapter }{LiftToSL2ZOdd}}
\logpage{[ 3, 3, 11 ]}\nobreak
\hyperdef{L}{X83FBC7C57B8CFB80}{}
{\noindent\textcolor{FuncColor}{$\triangleright$\enspace\texttt{LiftToSL2ZOdd({\mdseries\slshape G})\index{LiftToSL2ZOdd@\texttt{LiftToSL2ZOdd}}
\label{LiftToSL2ZOdd}
}\hfill{\scriptsize (operation)}}\\
\textbf{\indent Returns:\ }
A modular subgroup.



 Lifts a given subgroup \mbox{\texttt{\mdseries\slshape G}} of ${\rm PSL}_2({\ensuremath{\mathbb Z}})$ to an odd subgroup of ${\rm SL}_2({\ensuremath{\mathbb Z}})$, i.e. a group that does not contain $-1$ and whose projection to ${\rm PSL}_2({\ensuremath{\mathbb Z}})$ is \mbox{\texttt{\mdseries\slshape G}}. }

 
\begin{Verbatim}[commandchars=!@|,fontsize=\small,frame=single,label=Example]
  !gapprompt@gap>| !gapinput@G := ProjectiveModularSubgroup(|
  !gapprompt@>| !gapinput@(1,2)(3,7)(4,8)(5,9)(6,10)(11,12),|
  !gapprompt@>| !gapinput@(1,3,4)(2,5,6)(7,10,11)(8,12,9)|
  !gapprompt@>| !gapinput@);|
  <projective modular subgroup of index 12>
  !gapprompt@gap>| !gapinput@LiftToSL2ZOdd(G);|
  <modular subgroup of index 24>
\end{Verbatim}
 

\subsection{\textcolor{Chapter }{IndexModN (IndexModNPSL2Z)}}
\logpage{[ 3, 3, 12 ]}\nobreak
\hyperdef{L}{X81D105548473B68F}{}
{\noindent\textcolor{FuncColor}{$\triangleright$\enspace\texttt{IndexModN({\mdseries\slshape G, N})\index{IndexModN@\texttt{IndexModN}!IndexModNPSL2Z}
\label{IndexModN:IndexModNPSL2Z}
}\hfill{\scriptsize (operation)}}\\
\textbf{\indent Returns:\ }
A natural number.



 For a projective modular subgroup \mbox{\texttt{\mdseries\slshape G}} and a natural number \mbox{\texttt{\mdseries\slshape N}} this method calculates the index of the projection $\bar{G}$ of $G$ in ${\rm PSL}_2({\ensuremath{\mathbb Z}}/N{\ensuremath{\mathbb Z}})$. }

 
\begin{Verbatim}[commandchars=!@|,fontsize=\small,frame=single,label=Example]
  !gapprompt@gap>| !gapinput@G := ProjectiveModularSubgroup(|
  !gapprompt@>| !gapinput@(1,2)(3,6)(4,8)(5,9)(7,11)(10,13)(12,15)(14,17)(16,19)(18,21)(20,23)(22,24),|
  !gapprompt@>| !gapinput@(1,3,7,4)(2,5)(6,9,8,12,14,10)(11,13,16,20,18,15)(17,21,22,19)(23,24)|
  !gapprompt@>| !gapinput@);|
  <projective modular subgroup of index 24>
  !gapprompt@gap>| !gapinput@IndexModN(G, 2);|
  6
\end{Verbatim}
 

\subsection{\textcolor{Chapter }{Deficiency (DeficiencyPSL2Z)}}
\logpage{[ 3, 3, 13 ]}\nobreak
\hyperdef{L}{X858D7C9A7D134684}{}
{\noindent\textcolor{FuncColor}{$\triangleright$\enspace\texttt{Deficiency({\mdseries\slshape G, N})\index{Deficiency@\texttt{Deficiency}!DeficiencyPSL2Z}
\label{Deficiency:DeficiencyPSL2Z}
}\hfill{\scriptsize (operation)}}\\
\textbf{\indent Returns:\ }
A natural number.



 For a projective modular subgroup \mbox{\texttt{\mdseries\slshape G}} and a natural number \mbox{\texttt{\mdseries\slshape N}} this method calculates the so\texttt{\symbol{45}}called \emph{deficiency} of \mbox{\texttt{\mdseries\slshape G}} from being a congruence subgroup of level \mbox{\texttt{\mdseries\slshape N}}. \\
 The deficiency of a finite\texttt{\symbol{45}}index subgroup $\Gamma$ of ${\rm PSL}_2({\ensuremath{\mathbb Z}})$ was introduced in \cite{weitze_schmithuesen_deficiency}. It is defined as the index $[\Gamma(N) \colon \Gamma(N) \cap \Gamma]$ where $\Gamma(N)$ is the principal congruence subgroup of level $N$. }

 
\begin{Verbatim}[commandchars=!@|,fontsize=\small,frame=single,label=Example]
  !gapprompt@gap>| !gapinput@G := ProjectiveModularSubgroup(|
  !gapprompt@>| !gapinput@(1,2)(3,6)(4,8)(5,9)(7,11)(10,13)(12,15)(14,17)(16,19)(18,21)(20,23)(22,24),|
  !gapprompt@>| !gapinput@(1,3,7,4)(2,5)(6,9,8,12,14,10)(11,13,16,20,18,15)(17,21,22,19)(23,24)|
  !gapprompt@>| !gapinput@);|
  <projective modular subgroup of index 24>
  !gapprompt@gap>| !gapinput@Deficiency(G, 4);|
  4
\end{Verbatim}
 

\subsection{\textcolor{Chapter }{Deficiency (DeficiencyPSL2ZAttr)}}
\logpage{[ 3, 3, 14 ]}\nobreak
\hyperdef{L}{X78AE745D8277319E}{}
{\noindent\textcolor{FuncColor}{$\triangleright$\enspace\texttt{Deficiency({\mdseries\slshape G})\index{Deficiency@\texttt{Deficiency}!DeficiencyPSL2ZAttr}
\label{Deficiency:DeficiencyPSL2ZAttr}
}\hfill{\scriptsize (attribute)}}\\
\textbf{\indent Returns:\ }
A natural number.



 Shorthand for \texttt{Deficiency(G, GeneralizedLevel(G))}. }

 
\begin{Verbatim}[commandchars=!@|,fontsize=\small,frame=single,label=Example]
  !gapprompt@gap>| !gapinput@G := ProjectiveModularSubgroup(|
  !gapprompt@>| !gapinput@(1,2)(3,6)(4,8)(5,9)(7,11)(10,13)(12,15)(14,17)(16,19)(18,21)(20,23)(22,24),|
  !gapprompt@>| !gapinput@(1,3,7,4)(2,5)(6,9,8,12,14,10)(11,13,16,20,18,15)(17,21,22,19)(23,24)|
  !gapprompt@>| !gapinput@);|
  <projective modular subgroup of index 24>
  !gapprompt@gap>| !gapinput@Deficiency(G);|
  4
  !gapprompt@gap>| !gapinput@Deficiency(G, GeneralizedLevel(G));|
  4
\end{Verbatim}
 

\subsection{\textcolor{Chapter }{ConjugateGroup (ConjugateGroupPSL2Z)}}
\logpage{[ 3, 3, 15 ]}\nobreak
\hyperdef{L}{X82F9BB987AC90A9B}{}
{\noindent\textcolor{FuncColor}{$\triangleright$\enspace\texttt{ConjugateGroup({\mdseries\slshape G, A})\index{ConjugateGroup@\texttt{ConjugateGroup}!ConjugateGroupPSL2Z}
\label{ConjugateGroup:ConjugateGroupPSL2Z}
}\hfill{\scriptsize (operation)}}\\
\textbf{\indent Returns:\ }
A ProjectiveModularSubgroup.



 For an ${\rm SL}_2({\ensuremath{\mathbb Z}})$ matrix \mbox{\texttt{\mdseries\slshape A}}, conjugates the group \mbox{\texttt{\mdseries\slshape G}} by (the residue class in ${\rm PSL}_2({\ensuremath{\mathbb Z}})$ of) \mbox{\texttt{\mdseries\slshape A}} and returns the group $A^{-1}*G*A$. }

 

\subsection{\textcolor{Chapter }{NormalCore (NormalCorePSL2Z)}}
\logpage{[ 3, 3, 16 ]}\nobreak
\hyperdef{L}{X8772A1C282AC9172}{}
{\noindent\textcolor{FuncColor}{$\triangleright$\enspace\texttt{NormalCore({\mdseries\slshape G})\index{NormalCore@\texttt{NormalCore}!NormalCorePSL2Z}
\label{NormalCore:NormalCorePSL2Z}
}\hfill{\scriptsize (attribute)}}\\
\textbf{\indent Returns:\ }
A projective modular subgroup.



 Calculates the normal core of \mbox{\texttt{\mdseries\slshape G}} in ${\rm PSL}_2({\ensuremath{\mathbb Z}})$, i.e. the maximal subgroup of \mbox{\texttt{\mdseries\slshape G}} that is normal in ${\rm PSL}_2({\ensuremath{\mathbb Z}})$. }

 
\begin{Verbatim}[commandchars=!@|,fontsize=\small,frame=single,label=Example]
  !gapprompt@gap>| !gapinput@G := ProjectiveModularSubgroup(|
  !gapprompt@>| !gapinput@(1,2)(3,6)(4,8)(5,9)(7,11)(10,13)(12,15)(14,17)(16,19)(18,21)(20,23)(22,24),|
  !gapprompt@>| !gapinput@(1,3,7,4)(2,5)(6,9,8,12,14,10)(11,13,16,20,18,15)(17,21,22,19)(23,24)|
  !gapprompt@>| !gapinput@);|
  <projective modular subgroup of index 24>
  !gapprompt@gap>| !gapinput@NormalCore(G);|
  <projective modular subgroup of index 3456>
\end{Verbatim}
 

\subsection{\textcolor{Chapter }{QuotientByNormalCore (QuotientByNormalCorePSL2Z)}}
\logpage{[ 3, 3, 17 ]}\nobreak
\hyperdef{L}{X7F5807107867298D}{}
{\noindent\textcolor{FuncColor}{$\triangleright$\enspace\texttt{QuotientByNormalCore({\mdseries\slshape G})\index{QuotientByNormalCore@\texttt{QuotientByNormalCore}!QuotientByNormalCorePSL2Z}
\label{QuotientByNormalCore:QuotientByNormalCorePSL2Z}
}\hfill{\scriptsize (attribute)}}\\
\textbf{\indent Returns:\ }
A finite group.



 Calculates the quotient of ${\rm PSL}_2({\ensuremath{\mathbb Z}})$ by the normal core of \mbox{\texttt{\mdseries\slshape G}}. }

 
\begin{Verbatim}[commandchars=!@|,fontsize=\small,frame=single,label=Example]
  !gapprompt@gap>| !gapinput@G := ProjectiveModularSubgroup(|
  !gapprompt@>| !gapinput@(1,2)(3,6)(4,8)(5,9)(7,11)(10,13)(12,15)(14,17)(16,19)(18,21)(20,23)(22,24),|
  !gapprompt@>| !gapinput@(1,3,7,4)(2,5)(6,9,8,12,14,10)(11,13,16,20,18,15)(17,21,22,19)(23,24)|
  !gapprompt@>| !gapinput@);|
  <projective modular subgroup of index 24>
  !gapprompt@gap>| !gapinput@QuotientByNormalCore(G);|
  <permutation group with 2 generators>
\end{Verbatim}
 

\subsection{\textcolor{Chapter }{AssociatedCharacterTable (AssociatedCharacterTablePSL2Z)}}
\logpage{[ 3, 3, 18 ]}\nobreak
\hyperdef{L}{X7CDFA32A86529C5A}{}
{\noindent\textcolor{FuncColor}{$\triangleright$\enspace\texttt{AssociatedCharacterTable({\mdseries\slshape G})\index{AssociatedCharacterTable@\texttt{AssociatedCharacterTable}!AssociatedCharacterTablePSL2Z}
\label{AssociatedCharacterTable:AssociatedCharacterTablePSL2Z}
}\hfill{\scriptsize (attribute)}}\\
\textbf{\indent Returns:\ }
A character table.



 Returns the character table of ${\rm PSL}_2({\ensuremath{\mathbb Z}})/N$ where $N$ is the normal core of \mbox{\texttt{\mdseries\slshape G}}. }

 
\begin{Verbatim}[commandchars=!@|,fontsize=\small,frame=single,label=Example]
  !gapprompt@gap>| !gapinput@G := ProjectiveModularSubgroup(|
  !gapprompt@>| !gapinput@(1,2)(3,6)(4,8)(5,9)(7,11)(10,13)(12,15)(14,17)(16,19)(18,21)(20,23)(22,24),|
  !gapprompt@>| !gapinput@(1,3,7,4)(2,5)(6,9,8,12,14,10)(11,13,16,20,18,15)(17,21,22,19)(23,24)|
  !gapprompt@>| !gapinput@);|
  <projective modular subgroup of index 24>
  !gapprompt@gap>| !gapinput@AssociatedCharacterTable(G);|
  CharacterTable( <permutation group of size 3456 with 2 generators> )
\end{Verbatim}
 

\subsection{\textcolor{Chapter }{IsElementOf (IsElementOfPSL2Z)}}
\logpage{[ 3, 3, 19 ]}\nobreak
\hyperdef{L}{X7BF171DA7DBE5CDB}{}
{\noindent\textcolor{FuncColor}{$\triangleright$\enspace\texttt{IsElementOf({\mdseries\slshape A, G})\index{IsElementOf@\texttt{IsElementOf}!IsElementOfPSL2Z}
\label{IsElementOf:IsElementOfPSL2Z}
}\hfill{\scriptsize (operation)}}\\
\textbf{\indent Returns:\ }
True or false.



 This function checks if the image of a given matrix \mbox{\texttt{\mdseries\slshape A}} in ${\rm PSL}_2({\ensuremath{\mathbb Z}})$ is an element of the group \mbox{\texttt{\mdseries\slshape G}}. }

 
\begin{Verbatim}[commandchars=!@|,fontsize=\small,frame=single,label=Example]
  !gapprompt@gap>| !gapinput@G := ProjectiveModularSubgroup((1,2),(2,3));|
  <projective modular subgroup of index 3>
  !gapprompt@gap>| !gapinput@IsElementOf([[1,1],[0,1]], G);|
  true
  !gapprompt@gap>| !gapinput@IsElementOf([[0,-1],[1,0]], G);|
  false
\end{Verbatim}
 

\subsection{\textcolor{Chapter }{IsWordElementOf (IsWordElementOfPSL2ZString)}}
\logpage{[ 3, 3, 20 ]}\nobreak
\hyperdef{L}{X829E1DBD7A2AEE51}{}
{\noindent\textcolor{FuncColor}{$\triangleright$\enspace\texttt{IsWordElementOf({\mdseries\slshape s, G})\index{IsWordElementOf@\texttt{IsWordElementOf}!IsWordElementOfPSL2ZString}
\label{IsWordElementOf:IsWordElementOfPSL2ZString}
}\hfill{\scriptsize (operation)}}\\
\textbf{\indent Returns:\ }
True or false.



 This function checks if a given string \mbox{\texttt{\mdseries\slshape s}} represents an element of the group \mbox{\texttt{\mdseries\slshape G}}, written in the generators \mbox{\texttt{\mdseries\slshape S}} and \mbox{\texttt{\mdseries\slshape T}}. }

 
\begin{Verbatim}[commandchars=!@|,fontsize=\small,frame=single,label=Example]
  !gapprompt@gap>| !gapinput@G := ProjectiveModularSubgroup((1,2),(2,3));|
  <projective modular subgroup of index 3>
  !gapprompt@gap>| !gapinput@IsWordElementOf("S^2 * T", G);|
  true
\end{Verbatim}
 

\subsection{\textcolor{Chapter }{IsWordElementOf (IsWordElementOfPSL2Z)}}
\logpage{[ 3, 3, 21 ]}\nobreak
\hyperdef{L}{X7DDBB3A8848E83AE}{}
{\noindent\textcolor{FuncColor}{$\triangleright$\enspace\texttt{IsWordElementOf({\mdseries\slshape w, G})\index{IsWordElementOf@\texttt{IsWordElementOf}!IsWordElementOfPSL2Z}
\label{IsWordElementOf:IsWordElementOfPSL2Z}
}\hfill{\scriptsize (operation)}}\\
\textbf{\indent Returns:\ }
True or false.



 This function checks if a given word \mbox{\texttt{\mdseries\slshape w}} in the generators \mbox{\texttt{\mdseries\slshape S}} and \mbox{\texttt{\mdseries\slshape T}} represents an element of the group \mbox{\texttt{\mdseries\slshape G}}. }

 
\begin{Verbatim}[commandchars=!@|,fontsize=\small,frame=single,label=Example]
  !gapprompt@gap>| !gapinput@G := ProjectiveModularSubgroup((1,2),(2,3));|
  <projective modular subgroup of index 3>
  !gapprompt@gap>| !gapinput@F := FreeGroup("S", "T");|
  <free group on the generators [ S, T ]>
  !gapprompt@gap>| !gapinput@S := F.1;|
  S
  !gapprompt@gap>| !gapinput@T := F.2;|
  T
  !gapprompt@gap>| !gapinput@PSL2Z := F / ParseRelators([S, T], "S^2, (S*T)^3");|
  <fp group on the generators [ S, T ]>
  !gapprompt@gap>| !gapinput@S := GeneratorsOfGroup(PSL2Z)[1];|
  S
  !gapprompt@gap>| !gapinput@T := GeneratorsOfGroup(PSL2Z)[2];|
  T
  !gapprompt@gap>| !gapinput@IsWordElementOf(S^2 * T, G);|
  true
\end{Verbatim}
 

\subsection{\textcolor{Chapter }{Genus (GenusPSL2Z)}}
\logpage{[ 3, 3, 22 ]}\nobreak
\hyperdef{L}{X7A99E7BF7F5AEEC7}{}
{\noindent\textcolor{FuncColor}{$\triangleright$\enspace\texttt{Genus({\mdseries\slshape G})\index{Genus@\texttt{Genus}!GenusPSL2Z}
\label{Genus:GenusPSL2Z}
}\hfill{\scriptsize (attribute)}}\\
\textbf{\indent Returns:\ }
A non\texttt{\symbol{45}}negative integer.



 Computes the genus of the quotient $G \setminus {\ensuremath{\mathbb H}}$ via an algorithm described in \cite{schmithuesen_alg}. }

 
\begin{Verbatim}[commandchars=!@|,fontsize=\small,frame=single,label=Example]
  !gapprompt@gap>| !gapinput@G := ProjectiveModularSubgroup((1,2),(2,3));|
  <projective modular subgroup of index 3>
  !gapprompt@gap>| !gapinput@Genus(G);|
  0
\end{Verbatim}
 }

 }

 \def\bibname{References\logpage{[ "Bib", 0, 0 ]}
\hyperdef{L}{X7A6F98FD85F02BFE}{}
}

\bibliographystyle{alpha}
\bibliography{manual}

\addcontentsline{toc}{chapter}{References}

\def\indexname{Index\logpage{[ "Ind", 0, 0 ]}
\hyperdef{L}{X83A0356F839C696F}{}
}

\cleardoublepage
\phantomsection
\addcontentsline{toc}{chapter}{Index}


\printindex

\immediate\write\pagenrlog{["Ind", 0, 0], \arabic{page},}
\newpage
\immediate\write\pagenrlog{["End"], \arabic{page}];}
\immediate\closeout\pagenrlog
\end{document}
